% !TEX program = lualatex
% ---------------------------------------------------------------------------------------------
%  Readback note 7: Heights and the product formula; few translates nearly span; the height of one Gram determinant
%
%  Build:  python3 tools/rb.py note <slug>
%
%  One unit of a Lean artifact, rewritten as ordinary mathematics. The prose is hand-written. Lean names,
%  line ranges and counts come from data.tex; quoted Lean statements are copied from the artifact by the build.
%
%  Rules of the text: standard terminology; every symbol is explained where it first appears and listed again
%  in the notation at the end; no letter has two meanings; a proof is complete or says what it takes for granted;
%  Lean names appear only in the LEAN lines and in the concordance.
%
%  The signature rule (the build enforces it): a function, relation or predicate that the artifact defines is
%  written only through its notation macro, which the build makes from the spec (`notation`) with one slot for
%  each explicit argument of the Lean definition. No argument is fixed and dropped, also not "for brevity" and
%  not for one section. Every argument and every hypothesis of a stated declaration stands inside its
%  statement, or in a Setting that the statement cites by number; only implicit and instance arguments (the
%  ambient objects) may stand in the conventions.
% ---------------------------------------------------------------------------------------------
\documentclass[11pt,a4paper]{article}
\ifdefined\pdfvariable\pdfvariable suppressoptionalinfo \numexpr 1+2+4+8\relax\fi
\usepackage[a4paper,hmargin=35mm,top=22mm,bottom=27mm,footskip=12mm,nohead]{geometry}

\usepackage{fontspec}
\directlua{
  luaotfload.add_fallback("gfallback", {"STIXTwoMath-Regular.otf:mode=node;", "DejaVuSans.ttf:mode=node;"})
  luaotfload.add_fallback("gmonofallback", {"DejaVuSansMono.ttf:mode=node;", "STIXTwoMath-Regular.otf:mode=node;"})
}
\setmainfont{NewCM10}[Extension=.otf,UprightFont=*-Book,ItalicFont=*-BookItalic,
  BoldFont=*-Bold,BoldItalicFont=*-BoldItalic,SmallCapsFont=*-Book,SmallCapsFeatures={Letters=SmallCaps},
  RawFeature={fallback=gfallback}]
\setsansfont{NewCMSans10}[Extension=.otf,UprightFont=*-Book,ItalicFont=*-BookOblique,
  BoldFont=*-Bold,BoldItalicFont=*-BoldOblique,RawFeature={fallback=gfallback}]
\setmonofont{JetBrainsMono}[Extension=.otf,UprightFont=*-Regular,ItalicFont=*-Italic,
  BoldFont=*-SemiBold,BoldItalicFont=*-SemiBoldItalic,Scale=MatchLowercase,
  RawFeature={-calt;-liga;fallback=gmonofallback}]
\usepackage{amsmath}
\usepackage{amsthm}
\usepackage[math-style=TeX]{unicode-math}
\setmathfont{NewCMMath-Book.otf}
\emergencystretch=2.5em\relax

\usepackage{xcolor}
\usepackage{tikz}
\usetikzlibrary{calc,arrows.meta,positioning,backgrounds}
\usepackage[protrusion=true,expansion=true]{microtype}
\usepackage{enumitem}
\usepackage{needspace}
\usepackage{booktabs}
\usepackage{array}
\usepackage{longtable}
\usepackage{fancyhdr}
\usepackage{lastpage}
\usepackage{multicol}
\usepackage[hidelinks,pdfusetitle]{hyperref}

\definecolor{ink}{HTML}{0F172A}
\definecolor{inksoft}{HTML}{475569}
\definecolor{inkfaint}{HTML}{94A3B8}
\definecolor{rulegrey}{HTML}{CBD5E1}
\definecolor{arith}{HTML}{C2410C}
\definecolor{est}{HTML}{7C3AED}
\definecolor{ana}{HTML}{0891B2}
\definecolor{ok}{HTML}{047857}
\definecolor{leancol}{HTML}{3730A3}

\title{Heights and the product formula; few translates nearly span; the height of one Gram determinant}
\author{grh-compact}
\date{}

% ------------------------------------------------------------------------------------- macros
\newcommand{\leanfont}{\ttfamily\fontsize{8.6}{10.5}\selectfont\color{leancol}}
\newcommand{\lean}[1]{\csname LN@#1\endcsname}
\newcommand{\fillin}[1]{{\color{arith}\itshape[#1]}}  % a place of the template that is not written yet; the build reports every one that is left
\newcommand{\leanquote}[1]{\csname LQ@#1\endcsname}
% generated by rb.py note; do not edit
\newcommand{\ConcDecls}{27}
\newcommand{\Pin}{72fead10}
\newcommand{\PinText}{that of commit {\leanfont 72fead10}}
\newcommand{\ArtifactName}{grh-compact}
\expandafter\newcommand\csname LN@pairs\endcsname{\href{https://github.com/lsernau/grh-lean/blob/72fead102f663d5c7a407964a13dc0cf9d977207/grh-compact/src/New/Heights/Heights.lean\#L17-L17}{{\leanfont PrimePadicIso}}}
\expandafter\newcommand\csname LN@hle\endcsname{\href{https://github.com/lsernau/grh-lean/blob/72fead102f663d5c7a407964a13dc0cf9d977207/grh-compact/src/New/Heights/Heights.lean\#L21-L23}{{\leanfont HeightLe}}}
\expandafter\newcommand\csname LQ@hle\endcsname{\leanquoteline{0}{def HeightLe c (y : \_ → ℂ) := ∃ᵉ (H) (α) (e : ℕ → ℝ), H ≠ 0 ∧ (∀ σ : ℂ ≃+* ℂ, ‖σ H‖ ≤ α.exp) ∧}\leanquoteline{2}{(∀ ι : PrimePadicIso, ‖ι.2 (H * y ι)‖ ≤ (e ι.1).exp) ∧ ∀ᶠ M in atTop,}\leanquoteline{2}{α + ∑ l ∈ range M with l.Prime, e l ≤ c}}
\expandafter\newcommand\csname LN@pf\endcsname{\href{https://github.com/lsernau/grh-lean/blob/72fead102f663d5c7a407964a13dc0cf9d977207/grh-compact/src/New/Heights/Heights.lean\#L28-L30}{{\leanfont log\_product\_formula}}}
\expandafter\newcommand\csname LN@mono\endcsname{\href{https://github.com/lsernau/grh-lean/blob/72fead102f663d5c7a407964a13dc0cf9d977207/grh-compact/src/New/Heights/Heights.lean\#L36-L37}{{\leanfont HeightLe.\allowbreak{}mono}}}
\expandafter\newcommand\csname LN@ofle\endcsname{\href{https://github.com/lsernau/grh-lean/blob/72fead102f663d5c7a407964a13dc0cf9d977207/grh-compact/src/New/Heights/Heights.lean\#L42-L43}{{\leanfont HeightLe.\allowbreak{}of\_le}}}
\expandafter\newcommand\csname LN@mul\endcsname{\href{https://github.com/lsernau/grh-lean/blob/72fead102f663d5c7a407964a13dc0cf9d977207/grh-compact/src/New/Heights/Heights.lean\#L47-L50}{{\leanfont HeightLe.\allowbreak{}mul}}}
\expandafter\newcommand\csname LN@prod\endcsname{\href{https://github.com/lsernau/grh-lean/blob/72fead102f663d5c7a407964a13dc0cf9d977207/grh-compact/src/New/Heights/Heights.lean\#L54-L55}{{\leanfont HeightLe.\allowbreak{}prod}}}
\expandafter\newcommand\csname LN@pow\endcsname{\href{https://github.com/lsernau/grh-lean/blob/72fead102f663d5c7a407964a13dc0cf9d977207/grh-compact/src/New/Heights/Heights.lean\#L59-L60}{{\leanfont HeightLe.\allowbreak{}pow}}}
\expandafter\newcommand\csname LN@nonneg\endcsname{\href{https://github.com/lsernau/grh-lean/blob/72fead102f663d5c7a407964a13dc0cf9d977207/grh-compact/src/New/Heights/Heights.lean\#L65-L77}{{\leanfont HeightLe.\allowbreak{}nonneg}}}
\expandafter\newcommand\csname LQ@nonneg\endcsname{\leanquoteline{0}{lemma HeightLe.nonneg (h : HeightLe c 1) : 0 ≤ c}}
\expandafter\newcommand\csname LN@groupsum\endcsname{\href{https://github.com/lsernau/grh-lean/blob/72fead102f663d5c7a407964a13dc0cf9d977207/grh-compact/src/New/Heights/OneDeterminant.lean\#L22-L23}{{\leanfont inner\_sum\_sum\_eq\_card\_mul}}}
\expandafter\newcommand\csname LN@minors\endcsname{\href{https://github.com/lsernau/grh-lean/blob/72fead102f663d5c7a407964a13dc0cf9d977207/grh-compact/src/New/Heights/OneDeterminant.lean\#L28-L54}{{\leanfont card\_mul\_sum\_det\_le}}}
\expandafter\newcommand\csname LQ@minors\endcsname{\leanquoteline{0}{lemma card\_mul\_sum\_det\_le [Fintype k] (π : G →* GL n k) (A : \_ → Matrix n m k) \{Θ\} (h0 : ∀ s, 0 ≤ Θ}\leanquoteline{2}{s) (hu : ∀ a b, ⟪u a, u b⟫\_ℂ = Θ (a⁻¹ * b)) Y (hY : 0 ≤ Y) (h : ∀ b, ∀ v ≠ 0, ∑ s with A b *ᵥ v =}\leanquoteline{2}{cond b invTransposeGL 1 (π s) • A b *ᵥ v, Θ s ≤ Y) : card k * ∑ s with ((A true)ᵀ *}\leanquoteline{2}{(π s).1 * A false).det = 0, Θ s ≤ 2 * univ.sum Θ + card k \textasciicircum{} card m * √(card k \textasciicircum{} card n) * Y}}
\expandafter\newcommand\csname LN@dyad\endcsname{\href{https://github.com/lsernau/grh-lean/blob/72fead102f663d5c7a407964a13dc0cf9d977207/grh-compact/src/New/Heights/OneDeterminant.lean\#L58-L58}{{\leanfont dyadSum}}}
\expandafter\newcommand\csname LN@innerdyad\endcsname{\href{https://github.com/lsernau/grh-lean/blob/72fead102f663d5c7a407964a13dc0cf9d977207/grh-compact/src/New/Heights/OneDeterminant.lean\#L60-L62}{{\leanfont inner\_dyadSum}}}
\expandafter\newcommand\csname LN@variance\endcsname{\href{https://github.com/lsernau/grh-lean/blob/72fead102f663d5c7a407964a13dc0cf9d977207/grh-compact/src/New/Heights/OneDeterminant.lean\#L67-L82}{{\leanfont finrank\_sub\_le\_re\_inner\_dyadSum}}}
\expandafter\newcommand\csname LQ@variance\endcsname{\leanquoteline{0}{lemma finrank\_sub\_le\_re\_inner\_dyadSum [FiniteDimensional ℂ H] (P : EuclideanSpace ℂ n →ₗ[ℂ] H) (R :}\leanquoteline{2}{End ℂ H) (U A : Submodule ℂ \_) \{ε\} (k : ∑ x, ‖(P.adjoint ∘ₗ P - 1) (.single x 1)‖ ≤ ε)}\leanquoteline{2}{(h : ∀ v ∈ U, P v - R (P v) ∈ A) : finrank ℂ U - ε - ∑ x, ‖R (P (.single x 1))‖ \textasciicircum{} 2 ≤}\leanquoteline{2}{re ⟪dyadSum b (P <| .single · 1), dyadSum b (stdOrthonormalBasis ℂ A ·)⟫\_ℂ}}
\expandafter\newcommand\csname LN@spread\endcsname{\href{https://github.com/lsernau/grh-lean/blob/72fead102f663d5c7a407964a13dc0cf9d977207/grh-compact/src/New/Heights/OneDeterminant.lean\#L90-L92}{{\leanfont IsSpread}}}
\expandafter\newcommand\csname LQ@spread\endcsname{\leanquoteline{0}{def IsSpread \{W\} [AddCommMonoid W] [Module ℂ W] (B : S → V →ₗ[ℂ] W) η :=}\leanquoteline{2}{∀ F : Submodule .., let φ : ℝ := finrank ℂ F / finrank ℂ W}\leanquoteline{2}{∃ s, finrank ℂ (comap (B s) F) ≤ ((1 + φ) / 2 + η) * finrank ℂ V}}
\expandafter\newcommand\csname LN@few\endcsname{\href{https://github.com/lsernau/grh-lean/blob/72fead102f663d5c7a407964a13dc0cf9d977207/grh-compact/src/New/Heights/OneDeterminant.lean\#L97-L111}{{\leanfont exists\_few\_translates}}}
\expandafter\newcommand\csname LQ@few\endcsname{\leanquoteline{0}{lemma exists\_few\_translates \{B : Θ → V →ₗ[ℂ] W\} \{η\} (hη : 0 < η) (h : IsSpread B η) [Nontrivial V] :}\leanquoteline{4}{∃ t, \#t * finrank ℂ V ≤ 5 * finrank ℂ W * Real.log (1 + η⁻¹) ∧}\leanquoteline{4}{(1 - 4 * η) * finrank ℂ W ≤ finrank ℂ ↥(⨆ s ∈ t, (B s).range)}}
\expandafter\newcommand\csname LN@extend\endcsname{\href{https://github.com/lsernau/grh-lean/blob/72fead102f663d5c7a407964a13dc0cf9d977207/grh-compact/src/New/Heights/OneDeterminant.lean\#L115-L122}{{\leanfont exists\_extend\_basis\_of\_span\_eq\_top}}}
\expandafter\newcommand\csname LN@onedet\endcsname{\href{https://github.com/lsernau/grh-lean/blob/72fead102f663d5c7a407964a13dc0cf9d977207/grh-compact/src/New/Heights/OneDeterminant.lean\#L129-L149}{{\leanfont one\_determinant}}}
\expandafter\newcommand\csname LQ@onedet\endcsname{\leanquoteline{0}{lemma one\_determinant \{θ : \_ → Θ\} \{H : Θ → Type\} [∀ ϑ, NormedAddCommGroup (H ϑ)] [∀ ϑ,}\leanquoteline{2}{InnerProductSpace ℂ (H ϑ)] (u : ∀ ϑ, W →ₗ[ℂ] H ϑ) \{o\} (hu : (⇑(u o)).Injective) [Fintype S] \{x :}\leanquoteline{2}{S → W\} (i : LinearIndependent ℂ x) (hk : finrank ℂ W = card S) \{α δ g\} (hα : ∀ σ : ℂ ≃+* ℂ, ‖σ}\leanquoteline{2}{(gram ℂ (u o ∘ x)).det‖ ≤ Real.exp α) \{γ : \_ → ℂ\} (γ0 : ∀ ι, γ ι ≠ 0) (hγ : HeightLe g fun ι ↦ (‖γ}\leanquoteline{2}{ι‖ \textasciicircum{} 2 :)⁻¹) \{e : ℕ → ℝ\} (he : ∀ U, ∑ l ∈ U, e l ≤ δ) (Λ : ∀ ι : PrimePadicIso, (ℂ →+* ℂ\_[ι.1]) →}\leanquoteline{2}{Set (H (θ ι))) (k : ℕ) (w : ∀ ι j, ∀ x ∈ Λ ι (j.comp (starRingEnd ℂ)), ∀ y ∈ Λ ι j, ‖j ⟪x, y⟫\_ℂ‖ ≤}\leanquoteline{2}{1) (l : ∀ ι j, ∃ m, Real.exp (-e ι.1) * ‖j (γ ι)‖ \textasciicircum{} k ≤ ‖j (det m)‖ ∧ ∀ i, ∑ l, m i l • u \_ (x}\leanquoteline{2}{l) ∈ Λ ι j) : HeightLe (α + 2 * δ + k * g) fun ι ↦ (u (θ ι)).normDet \textasciicircum{} 2 / (u o).normDet \textasciicircum{} 2}}
\expandafter\newcommand\csname LN@algpadic\endcsname{\href{https://github.com/lsernau/grh-lean/blob/72fead102f663d5c7a407964a13dc0cf9d977207/grh-compact/src/Lib/CM/Genus/QuadraticZeta.lean\#L17-L24}{{\leanfont isAlgebraic\_of\_padic}}}
\expandafter\newcommand\csname LN@steinitz\endcsname{\href{https://github.com/lsernau/grh-lean/blob/72fead102f663d5c7a407964a13dc0cf9d977207/grh-compact/src/Lib/General/LinearAlgebra/CommonEigenvectors.lean\#L38-L47}{{\leanfont exists\_ringEquiv\_extend}}}
\expandafter\newcommand\csname LN@card\endcsname{\href{https://github.com/lsernau/grh-lean/blob/72fead102f663d5c7a407964a13dc0cf9d977207/grh-compact/src/Lib/PadicHodge/PeriodRings/Ainf.lean\#L16-L18}{{\leanfont mk\_complex\_eq\_mk\_padicAlgCl}}}
\expandafter\newcommand\csname LN@jiota\endcsname{\href{https://github.com/lsernau/grh-lean/blob/72fead102f663d5c7a407964a13dc0cf9d977207/grh-compact/src/Lib/ClassField/Places/UnramifiedCharacters.lean\#L44-L44}{{\leanfont complexToCp}}}
\expandafter\newcommand\csname LN@intmat\endcsname{\href{https://github.com/lsernau/grh-lean/blob/72fead102f663d5c7a407964a13dc0cf9d977207/grh-compact/src/Lib/General/LinearAlgebra/AdelicRingDefs.lean\#L53-L53}{{\leanfont IsIntMatrix}}}
\expandafter\newcommand\csname LN@detmem\endcsname{\href{https://github.com/lsernau/grh-lean/blob/72fead102f663d5c7a407964a13dc0cf9d977207/grh-compact/src/Lib/General/Algebra/GeneralisedIndex.lean\#L10-L11}{{\leanfont IsIntMatrix.\allowbreak{}det\_mem}}}
\expandafter\newcommand\csname LN@orth\endcsname{\href{https://github.com/lsernau/grh-lean/blob/72fead102f663d5c7a407964a13dc0cf9d977207/grh-compact/src/Lib/Weil/GaussSum/LocalHomUnit.lean\#L42-L46}{{\leanfont sum\_addChar\_dotProduct}}}
\expandafter\newcommand\csname LN@invT\endcsname{\href{https://github.com/lsernau/grh-lean/blob/72fead102f663d5c7a407964a13dc0cf9d977207/grh-compact/src/Lib/General/LinearAlgebra/ColumnLattices.lean\#L109-L110}{{\leanfont invTransposeGL}}}
\expandafter\newcommand\csname LN@transposeGL\endcsname{\href{https://github.com/lsernau/grh-lean/blob/72fead102f663d5c7a407964a13dc0cf9d977207/grh-compact/src/Lib/General/Groups/ValuationPair.lean\#L136-L136}{{\leanfont transposeGL}}}
\newcommand{\ConcordanceRows}{\crow{pairs}{\href{https://github.com/lsernau/grh-lean/blob/72fead102f663d5c7a407964a13dc0cf9d977207/grh-compact/src/New/Heights/Heights.lean\#L17-L17}{{\leanfont PrimePadicIso}}}{Heights.lean}{17--17}
\crow{hle}{\href{https://github.com/lsernau/grh-lean/blob/72fead102f663d5c7a407964a13dc0cf9d977207/grh-compact/src/New/Heights/Heights.lean\#L21-L23}{{\leanfont HeightLe}}}{Heights.lean}{21--23}
\crow{pf}{\href{https://github.com/lsernau/grh-lean/blob/72fead102f663d5c7a407964a13dc0cf9d977207/grh-compact/src/New/Heights/Heights.lean\#L28-L30}{{\leanfont log\_product\_formula}}}{Heights.lean}{28--30}
\crow{mono}{\href{https://github.com/lsernau/grh-lean/blob/72fead102f663d5c7a407964a13dc0cf9d977207/grh-compact/src/New/Heights/Heights.lean\#L36-L37}{{\leanfont HeightLe.\allowbreak{}mono}}}{Heights.lean}{36--37}
\crow{ofle}{\href{https://github.com/lsernau/grh-lean/blob/72fead102f663d5c7a407964a13dc0cf9d977207/grh-compact/src/New/Heights/Heights.lean\#L42-L43}{{\leanfont HeightLe.\allowbreak{}of\_le}}}{Heights.lean}{42--43}
\crow{mul}{\href{https://github.com/lsernau/grh-lean/blob/72fead102f663d5c7a407964a13dc0cf9d977207/grh-compact/src/New/Heights/Heights.lean\#L47-L50}{{\leanfont HeightLe.\allowbreak{}mul}}}{Heights.lean}{47--50}
\crow{prod}{\href{https://github.com/lsernau/grh-lean/blob/72fead102f663d5c7a407964a13dc0cf9d977207/grh-compact/src/New/Heights/Heights.lean\#L54-L55}{{\leanfont HeightLe.\allowbreak{}prod}}}{Heights.lean}{54--55}
\crow{pow}{\href{https://github.com/lsernau/grh-lean/blob/72fead102f663d5c7a407964a13dc0cf9d977207/grh-compact/src/New/Heights/Heights.lean\#L59-L60}{{\leanfont HeightLe.\allowbreak{}pow}}}{Heights.lean}{59--60}
\crow{nonneg}{\href{https://github.com/lsernau/grh-lean/blob/72fead102f663d5c7a407964a13dc0cf9d977207/grh-compact/src/New/Heights/Heights.lean\#L65-L77}{{\leanfont HeightLe.\allowbreak{}nonneg}}}{Heights.lean}{65--77}
\crow{groupsum}{\href{https://github.com/lsernau/grh-lean/blob/72fead102f663d5c7a407964a13dc0cf9d977207/grh-compact/src/New/Heights/OneDeterminant.lean\#L22-L23}{{\leanfont inner\_sum\_sum\_eq\_card\_mul}}}{OneDeterminant.lean}{22--23}
\crow{minors}{\href{https://github.com/lsernau/grh-lean/blob/72fead102f663d5c7a407964a13dc0cf9d977207/grh-compact/src/New/Heights/OneDeterminant.lean\#L28-L54}{{\leanfont card\_mul\_sum\_det\_le}}}{OneDeterminant.lean}{28--54}
\crow{dyad}{\href{https://github.com/lsernau/grh-lean/blob/72fead102f663d5c7a407964a13dc0cf9d977207/grh-compact/src/New/Heights/OneDeterminant.lean\#L58-L58}{{\leanfont dyadSum}}}{OneDeterminant.lean}{58--58}
\crow{innerdyad}{\href{https://github.com/lsernau/grh-lean/blob/72fead102f663d5c7a407964a13dc0cf9d977207/grh-compact/src/New/Heights/OneDeterminant.lean\#L60-L62}{{\leanfont inner\_dyadSum}}}{OneDeterminant.lean}{60--62}
\crow{variance}{\href{https://github.com/lsernau/grh-lean/blob/72fead102f663d5c7a407964a13dc0cf9d977207/grh-compact/src/New/Heights/OneDeterminant.lean\#L67-L82}{{\leanfont finrank\_sub\_le\_re\_inner\_dyadSum}}}{OneDeterminant.lean}{67--82}
\crow{spread}{\href{https://github.com/lsernau/grh-lean/blob/72fead102f663d5c7a407964a13dc0cf9d977207/grh-compact/src/New/Heights/OneDeterminant.lean\#L90-L92}{{\leanfont IsSpread}}}{OneDeterminant.lean}{90--92}
\crow{few}{\href{https://github.com/lsernau/grh-lean/blob/72fead102f663d5c7a407964a13dc0cf9d977207/grh-compact/src/New/Heights/OneDeterminant.lean\#L97-L111}{{\leanfont exists\_few\_translates}}}{OneDeterminant.lean}{97--111}
\crow{extend}{\href{https://github.com/lsernau/grh-lean/blob/72fead102f663d5c7a407964a13dc0cf9d977207/grh-compact/src/New/Heights/OneDeterminant.lean\#L115-L122}{{\leanfont exists\_extend\_basis\_of\_span\_eq\_top}}}{OneDeterminant.lean}{115--122}
\crow{onedet}{\href{https://github.com/lsernau/grh-lean/blob/72fead102f663d5c7a407964a13dc0cf9d977207/grh-compact/src/New/Heights/OneDeterminant.lean\#L129-L149}{{\leanfont one\_determinant}}}{OneDeterminant.lean}{129--149}
\crow{algpadic}{\href{https://github.com/lsernau/grh-lean/blob/72fead102f663d5c7a407964a13dc0cf9d977207/grh-compact/src/Lib/CM/Genus/QuadraticZeta.lean\#L17-L24}{{\leanfont isAlgebraic\_of\_padic}}}{QuadraticZeta.lean}{17--24}
\crow{steinitz}{\href{https://github.com/lsernau/grh-lean/blob/72fead102f663d5c7a407964a13dc0cf9d977207/grh-compact/src/Lib/General/LinearAlgebra/CommonEigenvectors.lean\#L38-L47}{{\leanfont exists\_ringEquiv\_extend}}}{CommonEigenvectors.lean}{38--47}
\crow{card}{\href{https://github.com/lsernau/grh-lean/blob/72fead102f663d5c7a407964a13dc0cf9d977207/grh-compact/src/Lib/PadicHodge/PeriodRings/Ainf.lean\#L16-L18}{{\leanfont mk\_complex\_eq\_mk\_padicAlgCl}}}{Ainf.lean}{16--18}
\crow{jiota}{\href{https://github.com/lsernau/grh-lean/blob/72fead102f663d5c7a407964a13dc0cf9d977207/grh-compact/src/Lib/ClassField/Places/UnramifiedCharacters.lean\#L44-L44}{{\leanfont complexToCp}}}{UnramifiedCharacters.lean}{44--44}
\crow{intmat}{\href{https://github.com/lsernau/grh-lean/blob/72fead102f663d5c7a407964a13dc0cf9d977207/grh-compact/src/Lib/General/LinearAlgebra/AdelicRingDefs.lean\#L53-L53}{{\leanfont IsIntMatrix}}}{AdelicRingDefs.lean}{53--53}
\crow{detmem}{\href{https://github.com/lsernau/grh-lean/blob/72fead102f663d5c7a407964a13dc0cf9d977207/grh-compact/src/Lib/General/Algebra/GeneralisedIndex.lean\#L10-L11}{{\leanfont IsIntMatrix.\allowbreak{}det\_mem}}}{GeneralisedIndex.lean}{10--11}
\crow{orth}{\href{https://github.com/lsernau/grh-lean/blob/72fead102f663d5c7a407964a13dc0cf9d977207/grh-compact/src/Lib/Weil/GaussSum/LocalHomUnit.lean\#L42-L46}{{\leanfont sum\_addChar\_dotProduct}}}{LocalHomUnit.lean}{42--46}
\crow{invT}{\href{https://github.com/lsernau/grh-lean/blob/72fead102f663d5c7a407964a13dc0cf9d977207/grh-compact/src/Lib/General/LinearAlgebra/ColumnLattices.lean\#L109-L110}{{\leanfont invTransposeGL}}}{ColumnLattices.lean}{109--110}
\crow{transposeGL}{\href{https://github.com/lsernau/grh-lean/blob/72fead102f663d5c7a407964a13dc0cf9d977207/grh-compact/src/Lib/General/Groups/ValuationPair.lean\#L136-L136}{{\leanfont transposeGL}}}{ValuationPair.lean}{136--136}}
% the notations of the defined functions: one slot for each explicit argument of the Lean declaration
\newcommand{\hle}[2]{\operatorname{ht}(#2)\le #1}
\newcommand{\dyad}[2]{\Delta_{#1}(#2)}
\newcommand{\spread}[2]{\operatorname{Spread}_{#2}(#1)}
\newcommand{\jof}[1]{j_{#1}}
\newcommand{\UnitMembers}{18}
\newcommand{\UnitInNote}{18}
\newcommand{\UnitProofs}{7}
\newcommand{\UnitSteps}{47}
\newcommand{\UnitAllSteps}{52}


\linespread{1.08}                       % 11 pt on 14.7 pt
\setlength{\parindent}{0pt}
\setlength{\parskip}{5.5pt plus 1.5pt minus 0.5pt}
\setlength{\abovedisplayskip}{7.5pt plus 2pt minus 3pt}
\setlength{\belowdisplayskip}{7.5pt plus 2pt minus 3pt}
\setlength{\abovedisplayshortskip}{3pt plus 2pt}
\setlength{\belowdisplayshortskip}{6pt plus 2pt minus 2pt}
\emergencystretch=1.5em
\clubpenalty=10000 \widowpenalty=10000 \displaywidowpenalty=10000
\relpenalty=7000 \binoppenalty=9000
\raggedbottom
\pagestyle{fancy}\fancyhf{}\renewcommand{\headrulewidth}{0pt}
\fancyfoot[C]{\sffamily\fontsize{7.6}{9}\selectfont\color{inkfaint}\thepage\,/\,\pageref*{LastPage}}
\setlist{topsep=2pt,partopsep=0pt,parsep=0pt,itemsep=4pt,leftmargin=2.1em,labelsep=0.55em}

\newcommand{\Q}{\mathbb{Q}}
\newcommand{\Z}{\mathbb{Z}}
\newcommand{\F}{\mathbb{F}}
\newcommand{\len}{\operatorname{len}}
\newcommand{\cok}{\operatorname{cok}}
\newcommand{\End}{\operatorname{End}}
\newcommand{\Hom}{\operatorname{Hom}}
\newcommand{\Gal}{\operatorname{Gal}}
\newcommand{\diag}{\operatorname{diag}}
\newcommand{\kw}[1]{\emph{#1}}
\newcommand{\m}[1]{\mbox{$#1$}}        % a short formula that must not be broken across lines
\newcommand{\head}[1]{\par\addvspace{17pt plus 4pt minus 2pt}\Needspace{6\baselineskip}%
  \noindent{\fontsize{13.2}{16}\selectfont\bfseries\color{ink}#1}\par\nopagebreak\vspace{-1pt}}
\newcommand{\sect}[2]{\refstepcounter{section}\label{#2}\par\addvspace{19pt plus 4pt minus 2pt}\Needspace{7\baselineskip}%
  \noindent\llap{{\fontsize{13.2}{16}\selectfont\bfseries\color{inkfaint}\thesection}\hspace{9pt}}%
  {\fontsize{13.2}{16}\selectfont\bfseries\color{ink}#1}\par\nopagebreak\vspace{-1pt}}
\newcommand{\subhead}[1]{\par\addvspace{9pt plus 2pt}\Needspace{4\baselineskip}\noindent\textbf{#1}\enspace\ignorespaces}
% where a statement lives in the Lean archive
\newcommand{\leantag}{{\sffamily\bfseries\fontsize{6.3}{8}\selectfont\color{leancol}\addfontfeatures{LetterSpace=7}LEAN}}
\newcommand{\sep}{\unskip\nobreak\hspace{5pt}{\color{inkfaint}\textperiodcentered}\hspace{5pt plus 2pt}\ignorespaces}
\newcommand{\leanline}[1]{\par\nopagebreak\vspace{-2.5pt}{\fontsize{8.9}{12.4}\selectfont\color{inksoft}\raggedright
  \noindent\llap{\leantag\hspace{9pt}}#1\par}\vspace{1pt}}
% a displayed statement in a tinted box: #1 colour, #2 text
\newcommand{\statement}[2]{\par\addvspace{7pt}\begingroup\setlength{\fboxsep}{8.5pt}\setlength{\fboxrule}{0.6pt}%
  \noindent\fcolorbox{#1!30!white}{#1!3!white}{\parbox{\dimexpr\linewidth-2\fboxsep-2\fboxrule\relax}{#2}}%
  \endgroup\par\addvspace{7pt}}
\newcommand{\nt}[2]{\par\noindent\hangindent=1.1em\hangafter=1 #1\enspace{\color{inksoft}#2}\par\vspace{1pt}}
\newcommand{\legdot}[1]{\tikz[baseline=-0.62ex]\fill[#1] (0,0) circle (2.5pt);\hspace{3pt}}
\newcommand{\crow}[4]{\csname CS@#1\endcsname & #2 & \textcolor{inksoft}{#3} & \textcolor{inksoft}{#4}\\}
\newcommand{\leanquoteline}[2]{\par\noindent\hangindent=\dimexpr#1\fontdimen2\font+2.2em\relax\hangafter=1
  \hspace*{#1\fontdimen2\font}#2}
\newcommand{\quoted}[1]{\par\vspace{3pt}{\ttfamily\fontsize{7.5}{10.4}\selectfont\color{leancol}\raggedright
  \setlength{\parskip}{0pt}\leanquote{#1}\par}\vspace{4pt}}
\newcommand{\bx}[3]{{\sffamily\bfseries\fontsize{7.3}{9.2}\selectfont\color{#1!82!black}#2}\\[1.5pt]#3}
\newcommand{\cs}[2]{\expandafter\newcommand\csname CS@#1\endcsname{#2}}

% theorems
\newtheoremstyle{plainx}{7pt plus 2pt}{7pt plus 2pt}{\itshape}{}{\bfseries}{.}{0.5em}{}
\newtheoremstyle{uprightx}{7pt plus 2pt}{7pt plus 2pt}{}{}{\bfseries}{.}{0.5em}{}
\theoremstyle{plainx}
\newtheorem{theorem}{Theorem}[section]
\newtheorem{lemma}[theorem]{Lemma}
\newtheorem{proposition}[theorem]{Proposition}
\newtheorem{corollary}[theorem]{Corollary}
\theoremstyle{uprightx}
\newtheorem{definition}[theorem]{Definition}
\newtheorem{setting}[theorem]{Setting}          % hypotheses shared by several statements; each of them cites it by number
\newtheorem{remark}[theorem]{Remark}
\newtheorem{example}[theorem]{Example}
\renewcommand{\qedsymbol}{{\setlength{\fboxsep}{0pt}\setlength{\fboxrule}{0.5pt}\fbox{\rule{0pt}{0.46em}\hspace{0.46em}}}}
% what holds throughout the note: the ambient objects only (implicit and instance arguments), never a hypothesis
\newenvironment{conventions}{\par\addvspace{5pt}\noindent\textbf{Conventions.}\enspace\ignorespaces}{\par\addvspace{3pt}}
\makeatletter
\renewenvironment{proof}[1][Proof]{\par\pushQED{\qed}\normalfont\topsep0pt\partopsep0pt
  \trivlist\item[\hskip\labelsep\itshape #1.]\ignorespaces}{\popQED\endtrivlist\@endpefalse}
\makeatother

% where each declaration of the concordance is treated: \cs{key}{place}, the place naming a \label or a lettered item
\cs{pairs}{Definition~\ref{d:pairs}}
\cs{hle}{Definition~\ref{d:height}}
\cs{pf}{Lemma~\ref{l:pf}}
\cs{mono}{Lemma~\ref{l:mono}}
\cs{ofle}{Lemma~\ref{l:ofle}}
\cs{mul}{Proposition~\ref{p:products}}
\cs{prod}{Proposition~\ref{p:products}}
\cs{pow}{Proposition~\ref{p:products}}
\cs{nonneg}{Theorem 1}
\cs{groupsum}{Lemma~\ref{l:groupsum}}
\cs{minors}{Proposition~\ref{p:minors}}
\cs{dyad}{Definition~\ref{d:dyad}}
\cs{innerdyad}{Lemma~\ref{l:innerdyad}}
\cs{variance}{Proposition~\ref{p:variance}}
\cs{spread}{Definition~\ref{d:spread}}
\cs{few}{Theorem 3}
\cs{extend}{Lemma~\ref{l:extend}}
\cs{onedet}{Theorem 2}
\cs{algpadic}{Input 1}
\cs{steinitz}{Input 2}
\cs{card}{Input 3}
\cs{jiota}{Definition~\ref{d:jiota}}
\cs{intmat}{Theorem 2}
\cs{detmem}{Theorem 2}
\cs{orth}{Proposition~\ref{p:minors}}
\cs{invT}{Proposition~\ref{p:minors}}
\cs{transposeGL}{Proposition~\ref{p:minors}}
\newcommand{\CC}{\mathbb{C}}
\newcommand{\RR}{\mathbb{R}}
\newcommand{\NN}{\mathbb{N}}
\newcommand{\QQ}{\mathbb{Q}}
\newcommand{\Ql}{\mathbb{Q}_\ell}
\newcommand{\Qlb}{\bar{\mathbb{Q}}_\ell}
\newcommand{\Cl}{\mathbb{C}_\ell}
\newcommand{\AutC}{\operatorname{Aut}(\mathbb{C})}
\newcommand{\IX}{\mathcal{I}}
\newcommand{\Hs}{\mathcal{H}}
\newcommand{\Ls}{\mathcal{L}}
\newcommand{\Ts}{\mathcal{T}}
\newcommand{\Us}{\mathcal{U}}
\newcommand{\re}{\operatorname{Re}}
\newcommand{\ord}{\operatorname{ord}}
\newcommand{\Nm}{\operatorname{Nm}}
\newcommand{\GL}{\operatorname{GL}}
\newcommand{\Gram}{\operatorname{Gram}}
\newcommand{\codim}{\operatorname{codim}}
\newcommand{\tp}{\mathrm{t}}
\newcommand{\one}{\mathbf{1}}
\newcommand{\abs}[1]{\lvert #1\rvert}
\newcommand{\norm}[1]{\lVert #1\rVert}
\newcommand{\ip}[2]{\langle #1,#2\rangle}
\newcommand{\bigip}[2]{\bigl\langle #1,#2\bigr\rangle}
\newcommand{\Row}{\mathrm{Row}}
\newcommand{\Col}{\mathrm{Col}}
\newcommand{\budget}{\mathrm{budget}}
\newcommand{\hmain}{h_{\mathrm{main}}}
\newcommand{\covol}{\operatorname{covol}}
\newcommand{\lf}[1]{{\leanfont #1}}

\begin{document}
\color{ink}

{\raggedright
{\sffamily\bfseries\fontsize{7.2}{9}\selectfont\color{leancol}\addfontfeatures{LetterSpace=8}READBACK NOTE 7 \textperiodcentered{} CHAPTER IV\par}
\vspace{5pt}
{\fontsize{21.5}{25.5}\selectfont\bfseries Heights and the product formula; few translates nearly span; the height of one Gram determinant\par}
\vspace{5pt}
{\fontsize{12.4}{16}\selectfont\itshape\color{inksoft} One unit of the Lean artifact \ArtifactName, read back into
mathematics\par}}
\vspace{8pt}
{\color{rulegrey}\hrule height 0.6pt}
\vspace{7pt}
\thispagestyle{empty}


\textbf{Summary.} A \emph{height bound} is defined for families $Z=(Z_{\ell,\iota})$ of complex numbers indexed by
the pairs of a prime~$\ell$ and a field isomorphism $\iota\colon\CC\to\Qlb$: one auxiliary number $H\ne0$ may be
multiplied in, and the bound adds the logarithmic size of all conjugates of~$H$ to the $\ell$-adic sizes of the
numbers $H\,Z_{\ell,\iota}$, summed over~$\ell$. Such bounds add in products, and the product formula, applied to
the norm of~$H$, shows that a bound for the constant family~$1$ is non-negative (Theorem~1). This lets the
argument compare $\ell$-adic estimates that are obtained with a different auxiliary object at every embedding.
Theorem~2 bounds the height of a family of ratios of Gram determinants, a twisted one at each embedding over one
untwisted one, by the size of the untwisted determinant, plus twice an index, plus $\nu$ times a height bound for
a family of scalars. The unit also holds linear algebra
for another part of the argument: if the translates of a linear map are spread, few of them nearly span the
target (Theorem~3); and two inequalities, one of Fourier type over a finite field, one of Hilbert--Schmidt type.

\statement{leancol}{\fontsize{9.6}{13.2}\selectfont\textbf{What this note is.} The artifact \ArtifactName{} is a
body of Lean code: the compact edition of a formal proof, due to L.~Alpöge with the AI system Claude, of the
Riemann hypothesis and of its generalisation to Dirichlet $L$-functions. That headline is a claim. The authors
report that every file compiles with Lean v4.35.0 against a pinned version of the library Mathlib and that Lean's
kernel accepted the proofs; an independent build is recorded in the handout's record of checks when it reports.
No mathematician has read the argument through, and nothing here says that the Riemann hypothesis is proved. This
note treats one unit of the artifact, the folder on heights: all \UnitInNote{} of its \UnitMembers{} declarations.
Their proofs are given in full. Taken for granted are three results of the background library of the artifact,
stated as Inputs~1 to~3 with the units that own them, and standard facts of Mathlib, which are named where they
are used. Every statement was compared with its Lean text, argument by argument. The Lean text read is \PinText;
Appendix~A is the concordance of \ConcDecls{} declarations.}

\sect{Heights, and the three main statements}{s:intro}
Throughout, $\ell$ is a prime number, $\Ql$ is the field of $\ell$-adic numbers, $\Qlb$ is an algebraic closure
of~$\Ql$, and $\Cl$ is the completion of~$\Qlb$. On these three fields $\abs{\,\cdot\,}_\ell$ is the absolute value
with \m{\abs{\ell}_\ell=1/\ell}, and $\abs{\,\cdot\,}$ is the usual absolute value of~$\CC$. The group $\AutC$
consists of all field automorphisms of~$\CC$, continuous or not.

\begin{definition}[the index set]\label{d:pairs}
$\IX$ is the set of all pairs $(\ell,\iota)$ of a prime number~$\ell$ and an isomorphism of rings
\m{\iota\colon\CC\to\Qlb}. In this note a \kw{family indexed by~$\IX$} is a family $Z=(Z_{\ell,\iota})$ of complex
numbers, one for each pair $(\ell,\iota)$ in~$\IX$. Such families are multiplied entry by entry, and $1$ also
denotes the family all of whose entries are~$1$.
\end{definition}
\leanline{\lean{pairs}}

Isomorphisms~$\iota$ exist for every~$\ell$, because $\CC$ and $\Qlb$ are algebraically closed fields of
characteristic~$0$ and of the same cardinality (Inputs 2 and~3 in §\ref{s:setting}). None of them is continuous
and none is distinguished. The index set is a way to speak of all embeddings into~$\Qlb$ of all number fields at
once: the restrictions of the isomorphisms~$\iota$ to a number field contained in~$\CC$ are exactly its field
embeddings into~$\Qlb$, again by Input~2. So for a complex number that is algebraic over~$\QQ$, its images
under all $\iota$ are its conjugates in~$\Qlb$, just as its images under $\AutC$ are its conjugates in~$\CC$.
The entries of a family may depend on $(\ell,\iota)$, and in the argument they do.

\begin{definition}[height at most $h$]\label{d:height}
Let $h$ be a real number and $Z=(Z_{\ell,\iota})$ a family indexed by~$\IX$. We write \m{\hle{h}{Z}}, and say that
$Z$ \kw{has height at most~$h$}, if there are a complex number \m{H\ne0}, a real number~$\alpha$ and a sequence
$\eta=(\eta_0,\eta_1,\eta_2,\dots)$ of real numbers such that
\begin{enumerate}[label=\textup{(\roman*)},leftmargin=2.3em,itemsep=2pt]
\item \m{\abs{\sigma(H)}\le\exp(\alpha)} for every $\sigma\in\AutC$;
\item \m{\abs{\iota(H\,Z_{\ell,\iota})}_\ell\le\exp(\eta_\ell)} for every $(\ell,\iota)\in\IX$;
\item there is a natural number~$n_0$ such that \m{\alpha+\sum_{\ell<n}\eta_\ell\le h} for every natural number
\m{n\ge n_0}, the sum being over the primes~$\ell$ below~$n$.
\end{enumerate}
\end{definition}
\leanline{\lean{hle}}

The sign \m{\hle{h}{Z}} is the name of a relation between $h$ and~$Z$; no number called the height of~$Z$ is defined
or needed. With \m{H=1} and \m{\alpha=0} the definition asks that the $\ell$-adic sizes of the entries, in
logarithms and summed over~$\ell$, be at most~$h$ (Lemma~\ref{l:ofle}). The number~$H$ allows more: it is one
number for all $(\ell,\iota)$, a common factor whose size at the archimedean places is paid for in~(i). The
sequence~$\eta$ is indexed by all natural numbers, as in the formal text; only its terms at the primes matter. Two
examples: for an integer \m{n\ge1}, the constant family~$1/n$ has height at most $\log n$ (take \m{H=n},
\m{\alpha=\log n}, \m{\eta=0}), and the constant family~$n$ has height at most $-\log n$ (take \m{H=1/n},
\m{\alpha=-\log n}, \m{\eta=0}). So a height bound says how small a family is $\ell$-adically, all primes together,
and it may be negative. It is not the height of Diophantine geometry: no positive parts are taken, the archimedean
places enter only through the auxiliary number~$H$, and the entries of a family need not be the conjugates of one
number. How negative a bound can be is limited by the first main result.

\statement{arith}{\phantomsection\label{t:1}\textbf{Theorem 1} (the product formula).\enspace\emph{Let $h$ be a
real number. If the constant family~$1$ has height at most~$h$, that is, if \m{\hle{h}{1}}, then \m{0\le h}.}}
\leanline{\lean{nonneg}}

This is the only place of the unit where arithmetic enters. The proof (§\ref{s:heights}) shows that the
number~$H$ of the definition is algebraic and applies the product formula of~$\QQ$ to its norm~$q$, a non-zero
rational number: the logarithms of the absolute values of~$q$ at all places of~$\QQ$ add up to zero, condition~(i)
bounds the logarithm at the archimedean place by $\alpha$ times the degree of~$H$, and condition~(ii) bounds the
logarithm at~$\ell$ by $\eta_\ell$ times the degree. So \m{\alpha+\sum_{\ell<n}\eta_\ell} cannot be negative once $n$
exceeds the primes that divide the numerator or the denominator of~$q$. The argument uses Theorem~1 together with the additivity of height bounds in products: if two families $Z$ and~$Z'$ have
entrywise product~$1$, then a bound for~$Z$ and a bound for~$Z'$ add up to a non-negative number
(Corollary~\ref{c:reciprocal}). An upper bound for the height of~$Z'$ thus becomes a lower bound for every height
bound of~$Z$, and this is how a covolume is bounded from below in the proof of Theorem~C of the argument
(§\ref{s:use}).

\subhead{Gram determinants.} Inner products are conjugate-linear in the first variable and linear in the second,
as in the library. For a finite family $v=(v_i)_{i\in I}$ of vectors of a complex inner product space,
\m{\Gram(v)=(\ip{v_i}{v_{i'}})_{i,i'\in I}} is its Gram matrix, and for a linear map~$T$ defined on these vectors
$Tv$ is the family $(Tv_i)_{i\in I}$. For a $\CC$-linear map $T\colon W\to\Hs$ from a finite-dimensional complex inner
product space to a complex inner product space, $T^{*}T$ is the endomorphism of~$W$ whose matrix in an orthonormal
basis~$v$ of~$W$ is $\Gram(Tv)$; it does not depend on~$v$. Its determinant $\det(T^{*}T)$ is a non-negative real
number, the square of the factor by which $T$ multiplies volumes, and it is non-zero exactly when $T$ is injective
(Lemma~\ref{l:gram}).

\begin{setting}[a family of linear maps, a basis, and integral structures]\label{set:gram}
Let $W$ be a finite-dimensional complex inner product space, $\Xi$ a set, and $\theta=(\theta_{\ell,\iota})$ a
family of elements of~$\Xi$, one for each pair $(\ell,\iota)$ in~$\IX$. For every $\vartheta\in\Xi$ let
$\Hs_\vartheta$ be a complex inner product space and \m{T_\vartheta\colon W\to\Hs_\vartheta} a $\CC$-linear map,
and let $\vartheta_0$ be an element of~$\Xi$. Let $I$ be a finite set and $v'=(v'_i)_{i\in I}$ a family of vectors
of~$W$. Let $\alpha$, $\delta$ and~$h$ be real numbers, $c=(c_{\ell,\iota})$ a family of complex numbers indexed
by~$\IX$, $\eta=(\eta_0,\eta_1,\eta_2,\dots)$ a sequence of real numbers, and $\nu$ a natural number. For every
$(\ell,\iota)\in\IX$ and every homomorphism of rings \m{j\colon\CC\to\Cl} let $\Ls_{\ell,\iota}(j)$ be a subset of
$\Hs_{\theta_{\ell,\iota}}$. For such~$j$, the homomorphism \m{\bar j\colon\CC\to\Cl} is $j$ preceded by complex
conjugation.
\end{setting}

\statement{arith}{\phantomsection\label{t:2}\textbf{Theorem 2} (the height of one Gram determinant).\enspace\emph{In
Setting~\ref{set:gram}, suppose that
\begin{enumerate}[label=\textup{(\alph*)},leftmargin=2.2em,itemsep=2pt,topsep=3pt]
\item the map $T_{\vartheta_0}$ is injective;
\item the family~$v'$ is linearly independent and \m{\dim W=\#I}, so that $v'$ is a basis of~$W$;
\item \m{\abs{\sigma(\det\Gram(T_{\vartheta_0}v'))}\le\exp(\alpha)} for every $\sigma\in\AutC$;
\item \m{c_{\ell,\iota}\ne0} for every $(\ell,\iota)\in\IX$, and \m{\hle{h}{\abs{c}^{-2}}} for the family
\m{\abs{c}^{-2}=(\abs{c_{\ell,\iota}}^{-2})};
\item \m{\sum_{n\in J}\eta_n\le\delta} for every finite set~$J$ of natural numbers;
\item for every $(\ell,\iota)\in\IX$, every homomorphism of rings \m{j\colon\CC\to\Cl}, every
$F\in\Ls_{\ell,\iota}(\bar j)$ and every $F'\in\Ls_{\ell,\iota}(j)$, one has \m{\abs{j(\ip{F}{F'})}_\ell\le1};
\item for every $(\ell,\iota)\in\IX$ and every homomorphism of rings \m{j\colon\CC\to\Cl} there is a matrix
\m{u=(u_{ii'})_{i,i'\in I}} of complex numbers such that
\[\exp(-\eta_\ell)\cdot\abs{j(c_{\ell,\iota})}_\ell^{\,\nu}\ \le\ \abs{j(\det u)}_\ell\]
and
\[\sum_{i'\in I}u_{ii'}\,T_{\theta_{\ell,\iota}}(v'_{i'})\in\Ls_{\ell,\iota}(j)\quad\text{for all } i\in I.\]
\end{enumerate}
Then the family $Z=(Z_{\ell,\iota})$ of the ratios
\[Z_{\ell,\iota}\ =\ \frac{\det(T_{\theta_{\ell,\iota}}^{*}T_{\theta_{\ell,\iota}})}{\det(T_{\vartheta_0}^{*}T_{\vartheta_0})}\]
satisfies \m{\hle{\alpha+2\delta+\nu h}{Z}}.}}
\leanline{\lean{onedet}}

In words: when $T_\vartheta$ is injective, the vectors $T_\vartheta v'_i$ form a basis of $T_\vartheta(W)$, and
$\det\Gram(T_\vartheta v')$ is the square of the volume of the parallelotope that they span (the squared covolume of
the overview of the handout, at the identity of~$\AutC$). Hypothesis~(c) bounds this number at the base index~$\vartheta_0$, at all its conjugates.
The sets $\Ls_{\ell,\iota}(j)$ play the part of the $j$-integral vectors: by~(f) the pairing of a $\bar j$-integral
vector with a $j$-integral one is $j$-integral. Hypothesis~(g) says that, at the index $\theta_{\ell,\iota}$ chosen
for the pair $(\ell,\iota)$, suitable combinations of the vectors $T_{\theta_{\ell,\iota}}v'_i$ are integral, with a
matrix of combination whose determinant is not too small: $\eta_\ell$ bounds the logarithm of an index, and
$c_{\ell,\iota}$ is a scalar that $\nu$ of the vectors are allowed to carry. The conclusion compares the covolume at
$\theta_{\ell,\iota}$ with the covolume at~$\vartheta_0$: the ratio of their squares is $\ell$-adically small, all
primes together, up to the archimedean size~$\alpha$, twice the index~$\delta$ and the cost~$h$ of each of the $\nu$
scalars. The index counts twice because the proof needs one integral combination for the embedding $\jof{\iota}$
of Definition~\ref{d:jiota} and one for its conjugate. The square of an index appears as it does for lattices over~$\Z$, where a lattice with integral Gram
matrix that has index~$n$ in another one forces $n^2$ times the Gram determinant of the latter to be an integer.

\subhead{Spread families.} The third statement is linear algebra over~$\CC$.

\begin{definition}[spread families of linear maps]\label{d:spread}
Let $U$ and $W$ be complex vector spaces, $\Omega$ a set, $B=(B_\gamma)_{\gamma\in\Omega}$ a family of $\CC$-linear
maps \m{B_\gamma\colon U\to W}, and $\varepsilon$ a real number. The family~$B$ is \kw{$\varepsilon$-spread}, in
symbols \m{\spread{B}{\varepsilon}}, if for every linear subspace \m{W'\subseteq W} there is a $\gamma\in\Omega$ with
\[\dim B_\gamma^{-1}(W')\ \le\ \Bigl(\frac12\Bigl(1+\frac{\dim W'}{\dim W}\Bigr)+\varepsilon\Bigr)\cdot\dim U .\]
\end{definition}
\leanline{\lean{spread}}

If $B_\gamma$ is injective and its image is in general position with respect to~$W'$, the preimage of~$W'$ has
about the share \m{\dim W'/\dim W} of the dimension of~$U$. The definition asks less: for every $W'$ some member of
the family must pull $W'$ back to at most the mean of this share and~$1$, up to~$\varepsilon$.

\statement{arith}{\phantomsection\label{t:3}\textbf{Theorem 3} (few translates nearly span).\enspace\emph{Let $U$ be
a finite-dimensional complex vector space with \m{U\ne0}, let $W$ be a finite-dimensional complex inner product
space, $\Omega$ a set, $B=(B_\gamma)_{\gamma\in\Omega}$ a family of $\CC$-linear maps \m{B_\gamma\colon U\to W}, and
\m{\varepsilon>0} a real number. If \m{\spread{B}{\varepsilon}}, then there is a finite subset
\m{\Ts\subseteq\Omega} with
\[\#\Ts\cdot\dim U\ \le\ 5\,\dim W\cdot\log\bigl(1+\varepsilon^{-1}\bigr)\]
and
\[(1-4\varepsilon)\,\dim W\ \le\ \dim\sum_{\gamma\in\Ts}B_\gamma(U).\]}}
\leanline{\lean{few}}

Each image $B_\gamma(U)$ has dimension at most $\dim U$, so at least \m{(1-4\varepsilon)\dim W/\dim U} members are
needed to span a subspace of dimension \m{(1-4\varepsilon)\dim W}. The theorem says that this least number, times
\m{5\log(1+\varepsilon^{-1})/(1-4\varepsilon)}, suffices. In the argument the maps $B_\gamma$ are the translates of
one map by the elements~$\gamma$ of a finite group, whence the name.

\sect{The setting}{s:setting}
\begin{conventions}\phantomsection\label{c:conventions}
The natural numbers are \m{0,1,2,\dots}, and $\#I$ is the number of elements of a finite set~$I$. The letter $\nu$ is a natural number, never a
character of the ideles (Convention~0.6). In Sections~\ref{s:intro} to~\ref{s:moments} the letter $\vartheta$ is an element of an
abstract set~$\Xi$; in Section~\ref{s:use} it is a twist of the data (Convention~0.5). The letter $v'$ is a second family of
vectors (not a finite place, Convention~0.4), and $R$ is a real number, not the polar functions of the overview. The
letter~$\ell$ always denotes a prime number; a sum \m{\sum_{\ell<n}} runs over the primes below~$n$. Vector spaces
are over~$\CC$ unless a finite field is named. A \kw{complex inner product space} is a complex vector space with a
positive definite hermitian form $\ip{\,\cdot\,}{\,\cdot\,}$, conjugate-linear in the first variable, and
\m{\norm{F}^2=\ip{F}{F}}; completeness is not assumed. For a finite set~$I$, the space $\CC^I$ has the inner product
for which the standard basis $\one=(\one_i)_{i\in I}$ is orthonormal. The tensor product of two complex inner product
spaces carries the inner product with
\m{\ip{F_1\otimes F_2}{F_1'\otimes F_2'}=\ip{F_1}{F_1'}\ip{F_2}{F_2'}}. For a finite-dimensional subspace~$\Hs'$ of an inner
product space~$\Hs$, $p_{\Hs'}$ is the orthogonal projection of~$\Hs$ onto~$\Hs'$. The functions of the library are defined everywhere:
the logarithm of a real number~$q$ is \m{\log\abs q} for \m{q\ne0} and $0$ for \m{q=0}; a quotient with
denominator~$0$ is~$0$; and the dimension of a vector space that is not finite-dimensional is~$0$. These
conventions enter two statements only, Lemma~\ref{l:pf} for \m{q=0} and Definition~\ref{d:spread}, which the
formal text states without assuming finite dimension. The declarations of this unit take no ambient object that
is fixed once for the whole note: each statement below introduces the spaces, sets and maps it is about.
\end{conventions}

\subhead{Letters local to this note.} Sections~\ref{s:intro} to~\ref{s:moments} do not depend on the data of the
proof, and ten letters have there a meaning of their own, recorded in the notation list: $\xi$ and $\rho$ (vectors
with coordinates in a finite set), $\tau$ and $\tau_0$ (homomorphisms of rings), $c$ (a family of scalars), $q$
(a rational number), $\Psi$ (an additive character), $\omega$ (a function on a group), $\Gamma_\vartheta$ (a Gram
determinant) and $\Omega$ (an index set). In the overview of the handout and in Chapter~III these letters denote
Dirichlet characters, zeros, infinite places, a residue class, the rational number of Theorem~A, cyclotomic
polynomials, residue characters, groups of rational points and a period. None of those objects occurs in
these sections, and of the local letters only $c$ returns in Section~\ref{s:use}, in its role of Theorem~2.

\begin{definition}[the embedding attached to $\iota$]\label{d:jiota}
For a prime number~$\ell$ and an isomorphism of rings \m{\iota\colon\CC\to\Qlb}, the homomorphism of rings
\m{\jof{\iota}\colon\CC\to\Cl} is $\iota$ followed by the inclusion of~$\Qlb$ in~$\Cl$. So
\m{\abs{\jof{\iota}(H)}_\ell=\abs{\iota(H)}_\ell} for every complex number~$H$.
\end{definition}
\leanline{\lean{jiota}\sep a definition of the background library, in the unit \emph{Lib.ClassField.Places}}

\statement{inksoft}{\textbf{Inputs from the background library.} The three results below are taken for granted
in this note. They are theorems of the background library of the artifact, which proves them from Mathlib; the
units that own them are named, and none of these units has been read yet.

\vspace{4pt}\phantomsection\label{i:1}\textbf{Input 1} (a number bounded at all $2$-adic identifications is
algebraic; unit \emph{Lib.CM.Genus}).\enspace \emph{Let $H$ be a complex number and $Y$ a real number. If
\m{\abs{\jof{\iota}(H)}_2\le Y} for every isomorphism of rings \m{\iota\colon\CC\to\bar{\QQ}_2}, then $H$ is
algebraic over~$\QQ$.}

\vspace{3pt}\phantomsection\label{i:2}\textbf{Input 2} (embeddings extend to isomorphisms; unit
\emph{Lib.General.LinearAlgebra}).\enspace \emph{Let $L$ be an algebraically closed field of the same cardinality
as~$\CC$, and let $\tau_0$ and~$\tau$ be injective homomorphisms of rings from one countable commutative ring
to~$\CC$ and to~$L$, respectively. Then there is an isomorphism of rings \m{\iota\colon\CC\to L} with
\m{\iota\circ\tau_0=\tau}.}

\vspace{3pt}\phantomsection\label{i:3}\textbf{Input 3} (cardinality; unit
\emph{Lib.PadicHodge.PeriodRings}).\enspace \emph{For every prime number~$\ell$, the field $\Qlb$ has the same
cardinality as~$\CC$.}}
\leanline{\lean{algpadic}\sep \lean{steinitz}\sep \lean{card}}

Input~2 is a form of Steinitz's theorem that an algebraically closed field is determined by its characteristic
and its transcendence degree: the catalogue of this handout records it under Steinitz~[St] as a special
case, that of two algebraically closed fields of characteristic~$0$ and of the cardinality of the continuum, with
an isomorphism that extends a given one between countable subrings; the place in~[St] is not confirmed.
For Inputs 1 and~3 the catalogue has no reference: reference to be confirmed. Both are elementary, the first
once Input~2 is granted. Input~3 holds because $\Ql$ has the cardinality of the continuum and an algebraic extension of an infinite
field has the cardinality of that field. Input~1 has a short formal proof, which was read: if $H$ were
transcendental, one would choose an element of~$\bar{\QQ}_2$ that is transcendental over~$\QQ$ and of absolute
value greater than~$Y$ (the algebraic elements are countable, so the others are dense), and Input~2, applied to
the two embeddings of the polynomial ring over~$\QQ$ in one variable that send the variable to~$H$ and to this
element, would give an~$\iota$ violating the bound.

\subhead{Facts of Mathlib.} The proofs below also use standard facts that the formal proofs take from Mathlib, the
mathematical library of Lean: unique factorisation of integers; that the image of the norm of a finite separable
field extension in an algebraically closed field is the product over the embeddings; that $\abs{\,\cdot\,}_\ell$
on~$\Qlb$ restricts to the $\ell$-adic absolute value on~$\QQ$; the properties of Gram matrices collected in
Lemma~\ref{l:gram}, where a short proof is given; Bessel's inequality, Parseval's identity and Pythagoras' theorem
in inner product spaces; the trace of an endomorphism as the sum of its diagonal coefficients in an orthonormal
basis; the existence of a non-trivial additive character of a finite field; and elementary linear algebra. Their
formal statements were not read: Mathlib is not part of the artifact, and its pinned version was not at hand when
this note was written. The dossier of the unit records these facts from their names and from the way they are
used.

\sect{Height bounds and the product formula}{s:heights}
In this section $Z$ and $Z'$ are families indexed by~$\IX$ and $h$, $h'$ are real numbers; a triple
$(H,\alpha,\eta)$ as in Definition~\ref{d:height} is called a \kw{witness} of \m{\hle{h}{Z}}.

\begin{lemma}[monotonicity]\label{l:mono}
Let $h$ and $h'$ be real numbers and $Z$, $Z'$ families indexed by~$\IX$. If \m{\hle{h}{Z}}, if \m{h\le h'}, and if
\m{\abs{\iota(Z'_{\ell,\iota})}_\ell\le\abs{\iota(Z_{\ell,\iota})}_\ell} for every $(\ell,\iota)\in\IX$, then
\m{\hle{h'}{Z'}}.
\end{lemma}
\begin{proof}
A witness $(H,\alpha,\eta)$ of \m{\hle{h}{Z}} is a witness of \m{\hle{h'}{Z'}}: conditions (i) and~(iii) are
unchanged or weaker, and
\[\abs{\iota(H Z'_{\ell,\iota})}_\ell=\abs{\iota(H)}_\ell\abs{\iota(Z'_{\ell,\iota})}_\ell\le\abs{\iota(H Z_{\ell,\iota})}_\ell\le\exp(\eta_\ell),\]
because $\iota$ is a homomorphism of rings and $\abs{\,\cdot\,}_\ell$ is multiplicative.
\end{proof}
\leanline{\lean{mono}}

\begin{lemma}[families of bounded entries]\label{l:ofle}
Let $h$ be a real number, $Z$ a family indexed by~$\IX$, and $\eta=(\eta_0,\eta_1,\eta_2,\dots)$ a sequence of real
numbers. If \m{\sum_{\ell\in J}\eta_\ell\le h} for every finite set~$J$ of prime numbers, and if
\m{\abs{\iota(Z_{\ell,\iota})}_\ell\le\exp(\eta_\ell)} for every $(\ell,\iota)\in\IX$, then \m{\hle{h}{Z}}.
\end{lemma}
\begin{proof}
The triple $(1,0,\eta)$ is a witness: \m{\sigma(1)=1} for every $\sigma\in\AutC$, condition~(ii) is the second
hypothesis, and (iii) holds for every~$n$ by the first hypothesis, applied to the set~$J$ of the primes below~$n$.
\end{proof}
\leanline{\lean{ofle}}

\begin{proposition}[height bounds add in products]\label{p:products}
Height bounds have the following three properties.
\begin{enumerate}[label=\textup{(\alph*)},leftmargin=2.2em,itemsep=2pt,topsep=2pt]
\item Let $h$ and $h'$ be real numbers and $Z$, $Z'$ families indexed by~$\IX$. If \m{\hle{h}{Z}} and
\m{\hle{h'}{Z'}}, then \m{\hle{h+h'}{ZZ'}}.
\item Let $\Omega$ be a set and $I\subseteq\Omega$ a finite subset, and for every $i\in\Omega$ let $h_i$ be a real
number and $Z_i$ a family indexed by~$\IX$. If \m{\hle{h_i}{Z_i}} for every $i\in I$, then
\m{\hle{\textstyle\sum_{i\in I}h_i}{\textstyle\prod_{i\in I}Z_i}}.
\item Let $h$ be a real number, $Z$ a family indexed by~$\IX$ with \m{\hle{h}{Z}}, and $n$ a natural number. Then
\m{\hle{n\,h}{Z^n}}.
\end{enumerate}
\end{proposition}
\begin{proof}
(a) Let $(H,\alpha,\eta)$ and $(H',\alpha',\eta')$ be witnesses, with thresholds $n_0$ and $n_0'$ in~(iii). Then
\m{HH'\ne0}, \m{\abs{\sigma(HH')}\le\exp(\alpha+\alpha')} for every~$\sigma$, and
\[\abs{\iota(HH'Z_{\ell,\iota}Z'_{\ell,\iota})}_\ell=\abs{\iota(HZ_{\ell,\iota})}_\ell\cdot\abs{\iota(H'Z'_{\ell,\iota})}_\ell\le\exp(\eta_\ell+\eta'_\ell)
\qquad\text{for every }(\ell,\iota)\in\IX .\]
For \m{n\ge n_0+n_0'} both inequalities~(iii) hold, and their sum is the inequality~(iii) for the triple
$(HH',\alpha+\alpha',\eta+\eta')$ and the bound \m{h+h'}.

(b) By induction on the number of elements of~$I$. If $I$ is empty, the product is the constant family~$1$ and the
sum is~$0$, and \m{\hle{0}{1}} by Lemma~\ref{l:ofle} with \m{\eta=0}. The step from $I$ to $I$ with one more element
is~(a).

(c) This is (b) for a set~$I$ of $n$ elements, with all $h_i$ equal to~$h$ and all $Z_i$ equal to~$Z$.
\end{proof}
\leanline{\lean{mul}\sep \lean{prod}\sep \lean{pow}}

\begin{lemma}[the product formula for $\QQ$]\label{l:pf}
Let $n$ be a natural number and $q$ a rational number, with numerator $q_{\mathrm{num}}$ and denominator
\m{q_{\mathrm{den}}\ge1} in lowest terms (for \m{q=0} these are $0$ and~$1$). If
\m{\max(\abs{q_{\mathrm{num}}},q_{\mathrm{den}})<n}, then
\[\log\abs q+\sum_{\ell<n}\log\abs q_\ell=0,\]
where $\abs q_\ell$ is the $\ell$-adic absolute value of~$q$, the sum is over the primes~$\ell$ below~$n$, and for
\m{q=0} all the logarithms are read as~$0$.
\end{lemma}
\begin{proof}
For \m{q=0} there is nothing to prove. Let \m{q\ne0}. A prime that divides $q_{\mathrm{num}}$ or $q_{\mathrm{den}}$
is at most \m{\max(\abs{q_{\mathrm{num}}},q_{\mathrm{den}})}, hence below~$n$. Write $\ord_\ell$ for the exponent
of~$\ell$ in an integer. By unique factorisation,
\[\log\abs{q_{\mathrm{num}}}=\sum_{\ell<n}\ord_\ell(q_{\mathrm{num}})\log\ell\qquad\text{and}\qquad
\log q_{\mathrm{den}}=\sum_{\ell<n}\ord_\ell(q_{\mathrm{den}})\log\ell .\]
Since \m{\log\abs q_\ell=(\ord_\ell(q_{\mathrm{den}})-\ord_\ell(q_{\mathrm{num}}))\log\ell}, the sum of the numbers
$\log\abs q_\ell$ over the primes \m{\ell<n} is \m{\log q_{\mathrm{den}}-\log\abs{q_{\mathrm{num}}}=-\log\abs q}.
\end{proof}
\leanline{\lean{pf}}

\begin{proof}[Proof of Theorem 1]
Let $(H,\alpha,\eta)$ be a witness of \m{\hle{h}{1}}, with threshold~$n_0$ in~(iii). Here (ii) reads
\m{\abs{\iota(H)}_\ell\le\exp(\eta_\ell)} for all $(\ell,\iota)\in\IX$.

\emph{The number $H$ is algebraic.} By (ii) for \m{\ell=2} and Definition~\ref{d:jiota},
\m{\abs{\jof{\iota}(H)}_2\le\exp(\eta_2)} for every isomorphism \m{\iota\colon\CC\to\bar{\QQ}_2}, and Input~1
applies. So $\QQ(H)$ is a number field. Write $\deg H$ for its degree over~$\QQ$, and let
\m{q=\Nm_{\QQ(H)/\QQ}(H)} be the norm of~$H$, a rational number, non-zero because \m{H\ne0}.

\emph{Local bounds.} We claim that
\begin{equation}\label{e:local}
\log\abs q\le\deg H\cdot\alpha\qquad\text{and}\qquad\log\abs q_\ell\le\deg H\cdot\eta_\ell\quad\text{for every
prime }\ell.
\end{equation}
Fix~$\ell$. The image of~$q$ in~$\Qlb$ is the product of the numbers $\tau(H)$ over the $\deg H$ field embeddings
\m{\tau\colon\QQ(H)\to\Qlb}. Each $\tau$ is the restriction of an isomorphism \m{\iota\colon\CC\to\Qlb}: this is
Input~2 for the field \m{L=\Qlb}, which has the cardinality of~$\CC$ by Input~3, for the countable ring $\QQ(H)$,
for its inclusion~$\tau_0$ in~$\CC$, and for~$\tau$. Hence \m{\abs{\tau(H)}_\ell=\abs{\iota(H)}_\ell\le\exp(\eta_\ell)}
by~(ii), and \m{\abs q_\ell\le\exp(\deg H\cdot\eta_\ell)}, because $\abs{\,\cdot\,}_\ell$ on~$\Qlb$ restricts to the
$\ell$-adic absolute value on~$\QQ$. The first inequality of~\eqref{e:local} is proved in the same way with $\CC$
in the place of~$\Qlb$: the embeddings \m{\QQ(H)\to\CC} are the restrictions of the automorphisms
$\sigma\in\AutC$ by Input~2 for \m{L=\CC}, and (i) bounds each $\abs{\sigma(H)}$ by $\exp(\alpha)$.

\emph{The product formula.} Put \m{n=n_0+\max(\abs{q_{\mathrm{num}}},q_{\mathrm{den}})+1}. By~(iii),
by~\eqref{e:local} and by Lemma~\ref{l:pf},
\[\deg H\cdot h\ \ge\ \deg H\cdot\Bigl(\alpha+\sum_{\ell<n}\eta_\ell\Bigr)\ \ge\ \log\abs q+\sum_{\ell<n}\log\abs q_\ell\ =\ 0,\]
and \m{\deg H\ge1}.
\end{proof}
\leanline{\lean{nonneg}\sep uses \lean{algpadic}, \lean{steinitz}, \lean{card} and \lean{pf}}

\begin{corollary}[the form in which Theorem 1 is used]\label{c:reciprocal}
Let $h$ and $h'$ be real numbers and $Z$, $Z'$ families indexed by~$\IX$ with \m{Z_{\ell,\iota}Z'_{\ell,\iota}=1}
for every $(\ell,\iota)\in\IX$. If \m{\hle{h}{Z}} and \m{\hle{h'}{Z'}}, then \m{h+h'\ge0}.
\end{corollary}
\begin{proof}
By Proposition~\ref{p:products}(a), \m{\hle{h+h'}{1}}; now apply Theorem~1.
\end{proof}
The corollary is not a declaration of the formal text; it is the last step of the formal proof of Theorem~C of
the argument, where it is applied to the square of one family and to a second family (§\ref{s:use}).

\begin{remark}[constant families]\label{r:constant}
The proof of Theorem~1 gives more, by the same steps. Let $Z$ be the constant family with value a non-zero
algebraic number, and let $\bar h$ be the mean of the logarithms of the absolute values of its conjugates
in~$\CC$. If \m{\hle{h}{Z}}, then \m{h\ge-\bar h}. Indeed, for a witness $(H,\alpha,\eta)$ the number $H$ is again
algebraic, and one applies the product formula to the norm of $H$ times the value of~$Z$, taken from the number
field that the two generate. So a height bound for a constant family cannot fall below what the product formula
allows, and the two examples of §\ref{s:intro} are best possible. This remark is ours; the formal text has only the
case of the constant family~$1$.
\end{remark}

\sect{The height of one Gram determinant}{s:gram}
\begin{lemma}[Gram matrices]\label{l:gram}
Let $I$ be a finite set and $F=(F_i)_{i\in I}$ a family of vectors of a complex inner product space.
\begin{enumerate}[label=\textup{(\arabic*)},leftmargin=2.2em,itemsep=2pt,topsep=2pt]
\item $\det\Gram(F)\ne0$ if and only if the family $F$ is linearly independent.
\item For matrices $u=(u_{ii'})$ and $u'=(u'_{ii'})$ of complex numbers indexed by \m{I\times I}, put
\m{F_i'=\sum_{i_1\in I}u'_{ii_1}F_{i_1}} and \m{F_i''=\sum_{i_1\in I}u_{ii_1}F_{i_1}} for $i\in I$. Then
\[\bigl(\ip{F_i'}{F_{i'}''}\bigr)_{i,i'\in I}\ =\ \overline{u'}\cdot\Gram(F)\cdot u^{\tp},\]
where $\overline{u'}$ is the matrix of the complex conjugates of the entries of~$u'$ and $u^{\tp}$ is the transpose
of~$u$.
\item Let $W$ be a finite-dimensional complex inner product space with \m{\dim W=\#I}, let $v'=(v'_i)_{i\in I}$ be a
family of vectors of~$W$, and let $T$ be a $\CC$-linear map from~$W$ to a complex inner product space. Then
\m{\det\Gram(Tv')=\det(T^{*}T)\cdot\det\Gram(v')}. In particular $\det(T^{*}T)\ne0$ if and only if $T$ is injective.
\end{enumerate}
\end{lemma}
\begin{proof}
Write $[F]$ for the linear map from $\CC^I$ that sends the standard basis vector~$\one_i$ to~$F_i$. Then $\Gram(F)$
is the matrix of the hermitian form \m{(\xi,\xi')\mapsto\ip{[F]\xi}{[F]\xi'}} on~$\CC^I$. This form is positive
semi-definite, and it is definite exactly when $[F]$ is injective, which is~(1). Statement~(2) is the expansion
of $\ip{F_i'}{F_{i'}''}$ by sesquilinearity. For~(3), choose an orthonormal basis~$v$ of~$W$ indexed by~$I$, and
let $u$ be the matrix with \m{v'_{i'}=\sum_{i\in I}u_{i'i}v_i}. By~(2), applied to the family~$Tv$ and then to the
family~$v$, \m{\Gram(Tv')=\overline{u}\cdot\Gram(Tv)\cdot u^{\tp}} and \m{\Gram(v')=\overline{u}\cdot u^{\tp}}. As
$\Gram(Tv)$ is the matrix of $T^{*}T$ in the basis~$v$, taking determinants gives the formula. The last assertion
follows from~(1) for the family~$Tv$.
\end{proof}
\leanline{facts of Mathlib; the formal text writes $\det(T^{*}T)$ as the square of a function of Mathlib that is
the volume factor of a linear map}

\begin{proof}[Proof of Theorem 2]
For every $\vartheta$ in~$\Xi$ put \m{\Gamma_\vartheta=\det\Gram(T_\vartheta v')}. By~(b) and
Lemma~\ref{l:gram}(1) the number $\det\Gram(v')$ is not zero, and by Lemma~\ref{l:gram}(3)
\[\Gamma_\vartheta\ =\ \det(T_\vartheta^{*}T_\vartheta)\cdot\det\Gram(v')\qquad\text{for every }\vartheta\in\Xi .\]
By~(a), \m{\Gamma_{\vartheta_0}\ne0}. So the family of the theorem is
\m{Z_{\ell,\iota}=\Gamma_{\theta_{\ell,\iota}}/\Gamma_{\vartheta_0}}.

\emph{Reduction.} Let $Z'$ be the family with entries \m{Z'_{\ell,\iota}=Z_{\ell,\iota}\,\abs{c_{\ell,\iota}}^{2\nu}}. Then
$Z$ is the product of~$Z'$ and of the $\nu$-th power of the family $\abs{c}^{-2}$, and this power has height at
most~$\nu h$ by~(d) and Proposition~\ref{p:products}(c). By Proposition~\ref{p:products}(a) it suffices to show
that \m{\hle{\alpha+2\delta}{Z'}}. We claim that the triple $(\Gamma_{\vartheta_0},\alpha,2\eta)$ is a witness.
Condition~(i) is~(c). Condition~(iii) holds for every~$n$, by~(e) for the set~$J$ of the primes below~$n$.
Condition~(ii) is the inequality
\begin{equation}\label{e:ladic}
\bigl\lvert\iota\bigl(\Gamma_{\theta_{\ell,\iota}}\cdot\abs{c_{\ell,\iota}}^{2\nu}\bigr)\bigr\rvert_\ell\ \le\ \exp(2\eta_\ell)
\qquad\text{for every }(\ell,\iota)\in\IX,
\end{equation}
which remains to be proved.

\emph{The $\ell$-adic bound.} Fix $(\ell,\iota)\in\IX$ and put \m{j=\jof{\iota}}, so that
\m{\abs{\iota(\,\cdot\,)}_\ell=\abs{j(\,\cdot\,)}_\ell}. Let $u$ be a matrix as in~(g) for $(\ell,\iota)$ and~$j$,
and $u'$ a matrix as in~(g) for $(\ell,\iota)$ and~$\bar j$. Since $\bar j$ is $j$ preceded by complex conjugation,
the two inequalities of~(g) read
\begin{equation}\label{e:index}
\exp(-\eta_\ell)\,\abs{j(c_{\ell,\iota})}_\ell^{\,\nu}\le\abs{j(\det u)}_\ell
\qquad\text{and}\qquad
\exp(-\eta_\ell)\,\abs{j(\overline{c_{\ell,\iota}})}_\ell^{\,\nu}\le\abs{j(\det\overline{u'})}_\ell .
\end{equation}
Put \m{F_i'=\sum_{i_1\in I}u'_{ii_1}T_{\theta_{\ell,\iota}}(v'_{i_1})} and
\m{F_i''=\sum_{i_1\in I}u_{ii_1}T_{\theta_{\ell,\iota}}(v'_{i_1})} for $i\in I$. By~(g) the vectors $F_i'$ lie in
$\Ls_{\ell,\iota}(\bar j)$ and the vectors $F_i''$ lie in $\Ls_{\ell,\iota}(j)$, so by~(f) every entry of the matrix
$(\ip{F_i'}{F_{i'}''})_{i,i'\in I}$ lies in the set~$O_j$ of the complex numbers whose image under~$j$ has absolute
value at most~$1$. This set is a subring of~$\CC$, because $\abs{\,\cdot\,}_\ell$ is non-archimedean; hence the
determinant of the matrix lies in~$O_j$ as well. By Lemma~\ref{l:gram}(2) the matrix is
\m{\overline{u'}\cdot\Gram(T_{\theta_{\ell,\iota}}v')\cdot u^{\tp}}, with determinant
\m{\det\overline{u'}\cdot\Gamma_{\theta_{\ell,\iota}}\cdot\det u}. Therefore
\begin{equation}\label{e:integral}
\abs{j(\det\overline{u'})}_\ell\cdot\abs{j(\Gamma_{\theta_{\ell,\iota}})}_\ell\cdot\abs{j(\det u)}_\ell\ \le\ 1 .
\end{equation}
Now \m{\abs{c_{\ell,\iota}}^{2}=\overline{c_{\ell,\iota}}\cdot c_{\ell,\iota}}, so by \eqref{e:index} and
then~\eqref{e:integral}
\[\bigl\lvert j\bigl(\Gamma_{\theta_{\ell,\iota}}\abs{c_{\ell,\iota}}^{2\nu}\bigr)\bigr\rvert_\ell
\ \le\ \abs{j(\Gamma_{\theta_{\ell,\iota}})}_\ell\cdot\exp(2\eta_\ell)\cdot\abs{j(\det\overline{u'})}_\ell\cdot\abs{j(\det u)}_\ell
\ \le\ \exp(2\eta_\ell).\]
This is~\eqref{e:ladic}.
\end{proof}
\leanline{\lean{onedet}\sep uses \lean{mul}, \lean{pow}, \lean{jiota}, and from the background library
\lean{intmat} and \lean{detmem}: a determinant of a matrix with entries in a subring lies in the subring}

The proof uses hypotheses (f) and~(g) only for the two homomorphisms $\jof{\iota}$ and its conjugate, and it uses
nothing about the sets $\Ls_{\ell,\iota}(j)$ beyond (f) and~(g). The inner product of~$W$ does not matter either:
by Lemma~\ref{l:gram}(3) the ratio $Z_{\ell,\iota}$ is a ratio of Gram determinants of images of the basis~$v'$.

\subhead{Completing a family to a basis.} Theorem~2 needs a basis of~$W$. In the argument a linearly independent
family is given, and it is completed by vectors of a second kind, for which hypothesis~(g) holds with the
scalars~$c_{\ell,\iota}$. The following lemma provides the completion.

\begin{lemma}[extension to a basis from a spanning set]\label{l:extend}
Let $W$ be a finite-dimensional complex vector space, $n$ a natural number, $S$ a subset of~$W$ that spans~$W$,
and \m{y_1,\dots,y_n} linearly independent vectors of~$W$. Then there are a natural number~$n'$ and vectors
\m{z_1,\dots,z_{n'}} in~$S$ such that \m{y_1,\dots,y_n,z_1,\dots,z_{n'}} are linearly independent and
\m{\dim W=n+n'}.
\end{lemma}
\begin{proof}
Let $W_y$ be the span of \m{y_1,\dots,y_n}; it has dimension~$n$. The images of the elements of~$S$ span the
quotient $W/W_y$, so there are \m{z_1,\dots,z_{n'}} in~$S$ whose images form a basis of $W/W_y$, and
\m{n'=\dim W-n}. A linear relation among \m{y_1,\dots,y_n,z_1,\dots,z_{n'}} becomes, modulo~$W_y$, a relation among
the images of the~$z_i$, so its coefficients at the~$z_i$ vanish; what is left is a relation among the~$y_i$, whose
coefficients vanish too.
\end{proof}
\leanline{\lean{extend}}

\sect{Few translates nearly span}{s:few}
\begin{proof}[Proof of Theorem 3]
If \m{\dim W=0}, the empty set has the two properties. Let \m{\dim W\ge1}. For a finite subset
\m{\Ts'\subseteq\Omega} put
\[W_{\Ts'}=\sum_{\gamma\in\Ts'}B_\gamma(U),\qquad\codim W_{\Ts'}=\dim W-\dim W_{\Ts'},\]
\[\Pi(\Ts')\ =\ \#\Ts'\cdot\dim U\ +\ 5\,\dim W\cdot\log\Bigl(1+\frac{\codim W_{\Ts'}}{\varepsilon\,\dim W}\Bigr).\]
Call $\Ts'$ \emph{admissible} if \m{\Pi(\Ts')\le5\,\dim W\cdot\log(1+\varepsilon^{-1})}. The empty set is admissible,
with equality. Among the admissible sets choose one, $\Ts$, for which $\dim W_{\Ts}$ is largest. The logarithm in
$\Pi(\Ts)$ is non-negative, so
\[\#\Ts\cdot\dim U\ \le\ \Pi(\Ts)\ \le\ 5\,\dim W\cdot\log(1+\varepsilon^{-1}):\]
this is the first inequality.

For the second, suppose that \m{\dim W_{\Ts}<(1-4\varepsilon)\dim W}, that is,
\begin{equation}\label{e:cd}
\codim W_{\Ts}\ >\ 4\,\varepsilon\,\dim W .
\end{equation}
Since $B$ is $\varepsilon$-spread, there is a $\gamma\in\Omega$ with
\[\dim B_\gamma^{-1}(W_{\Ts})\ \le\ \Bigl(\frac12\Bigl(1+\frac{\dim W_{\Ts}}{\dim W}\Bigr)+\varepsilon\Bigr)\dim U
\ =\ \Bigl(1-\frac{\codim W_{\Ts}}{2\dim W}+\varepsilon\Bigr)\dim U .\]
Let $\mathrm{rk}$ be the rank of the map \m{U\to W/W_{\Ts}} induced by~$B_\gamma$. Its kernel is
$B_\gamma^{-1}(W_{\Ts})$, so
\begin{equation}\label{e:rk}
\mathrm{rk}\ =\ \dim U-\dim B_\gamma^{-1}(W_{\Ts})\ \ge\ \Bigl(\frac{\codim W_{\Ts}}{2\dim W}-\varepsilon\Bigr)\dim U
\ >\ \frac{\codim W_{\Ts}}{4\dim W}\cdot\dim U\ >\ 0,
\end{equation}
the last two inequalities by~\eqref{e:cd} and \m{\dim U\ge1}. Put \m{\Ts''=\Ts\cup\{\gamma\}}. Then
\m{W_{\Ts''}=W_{\Ts}+B_\gamma(U)} and \m{\codim W_{\Ts''}=\codim W_{\Ts}-\mathrm{rk}}, so that the quotient of the
arguments of the two logarithms is \m{1-\mathrm{rk}/(\varepsilon\dim W+\codim W_{\Ts})}. This number lies in the
interval from $0$ to~$1$, and the logarithm of a positive number is at most that number minus~$1$. As $\Ts''$ has
at most one element more than~$\Ts$,
\begin{align*}
\Pi(\Ts'')-\Pi(\Ts)\ &\le\ \dim U+5\,\dim W\cdot\log\Bigl(1-\frac{\mathrm{rk}}{\varepsilon\dim W+\codim W_{\Ts}}\Bigr)\\
&\le\ \dim U-\frac{5\,\dim W\cdot\mathrm{rk}}{\varepsilon\dim W+\codim W_{\Ts}} .
\end{align*}
By the first inequality of~\eqref{e:rk} and by~\eqref{e:cd},
\[5\,\dim W\cdot\mathrm{rk}\ \ge\ \Bigl(\frac52\codim W_{\Ts}-5\,\varepsilon\dim W\Bigr)\dim U
\ \ge\ \bigl(\codim W_{\Ts}+\varepsilon\dim W\bigr)\dim U ,\]
the second step because \m{\tfrac32\codim W_{\Ts}\ge6\,\varepsilon\dim W}. Hence \m{\Pi(\Ts'')\le\Pi(\Ts)}, and
$\Ts''$ is admissible. By the choice of~$\Ts$, \m{\dim W_{\Ts''}\le\dim W_{\Ts}}, so \m{B_\gamma(U)\subseteq W_{\Ts}}
and \m{\mathrm{rk}=0}, against~\eqref{e:rk}.
\end{proof}
\leanline{\lean{few}\sep uses \lean{spread}}

The potential $\Pi$ explains the two constants. Adding a member costs $\dim U$ in the first term. The spread
property guarantees that a suitable member lowers the codimension by a fixed share of what is left, as long as
\eqref{e:cd} holds, and the logarithm turns this into a drop of the second term by at least $\dim U$. The
factor~$5$ is what makes the last display true under~\eqref{e:cd}, and $4\varepsilon$ is where the process stops.

\sect{Two inequalities behind the spread property}{s:moments}
The hypothesis of Theorem~3 is proved in another chapter, by a second moment over a finite group. The unit on
heights holds the two general inequalities of that proof. Both concern a finite group~$G$ and a map from~$G$ to an
inner product space whose inner products depend only on the quotient of the two arguments.

\begin{lemma}[inner products of sums over a group]\label{l:groupsum}
Let $G$ be a finite group, $\Hs$ a complex inner product space, \m{v\colon G\to\Hs} and \m{v'\colon G\to\Hs} two maps,
and \m{\omega\colon G\to\CC} a function such that \m{\ip{v(\gamma)}{v'(\gamma')}=\omega(\gamma^{-1}\gamma')} for all
$\gamma,\gamma'\in G$. Then
\[\bigip{\,\sum_{\gamma\in G}v(\gamma)}{\sum_{\gamma\in G}v'(\gamma)}\ =\ \#G\cdot\sum_{\gamma\in G}\omega(\gamma).\]
\end{lemma}
\begin{proof}
The left side is \m{\sum_{\gamma}\sum_{\gamma'}\omega(\gamma^{-1}\gamma')}, and for each~$\gamma$ the map
\m{\gamma'\mapsto\gamma^{-1}\gamma'} is a bijection of~$G$.
\end{proof}
\leanline{\lean{groupsum}}

In the next statement $\F$ is a finite field, $\F^{I}$ is the space of vectors indexed by a finite set~$I$, and
$\GL_I(\F)$ is the group of invertible matrices indexed by \m{I\times I}, which acts on~$\F^I$. A function $\omega$ as
in~(b) below is a function of positive type on~$G$.

\begin{proposition}[singular minors are rare]\label{p:minors}
Let $G$ be a finite group, $\Hs$ a complex inner product space and \m{v\colon G\to\Hs} a map. Let $\F$ be a finite
field, $I$ and $I'$ finite sets, \m{\pi\colon G\to\GL_I(\F)} a homomorphism of groups, and $A=(A_1,A_0)$ a pair of
matrices over~$\F$ with rows indexed by~$I$ and columns indexed by~$I'$. Let \m{\omega\colon G\to\RR} be a function
and $Y$ a real number. Suppose that
\begin{enumerate}[label=\textup{(\alph*)},leftmargin=2.2em,itemsep=2pt,topsep=2pt]
\item \m{\omega(\gamma)\ge0} for every $\gamma\in G$;
\item \m{\ip{v(\gamma)}{v(\gamma')}=\omega(\gamma^{-1}\gamma')} for all $\gamma,\gamma'\in G$;
\item \m{Y\ge0};
\item for every non-zero vector \m{\xi\in\F^{I'}}, the sum of $\omega(\gamma)$ over all $\gamma\in G$ such that
\m{(\pi(\gamma)^{\tp})^{-1}A_1\xi=A_1\xi} is at most~$Y$, and so is the sum of $\omega(\gamma)$ over all
$\gamma\in G$ such that \m{\pi(\gamma)A_0\xi=A_0\xi}.
\end{enumerate}
Then
\[\#\F\cdot\sum_{\substack{\gamma\in G\\ \det(A_1^{\tp}\pi(\gamma)A_0)=0}}\omega(\gamma)\ \ \le\ \ 2\sum_{\gamma\in G}\omega(\gamma)
\ +\ (\#\F)^{\#I'}\sqrt{(\#\F)^{\#I}}\cdot Y .\]
\end{proposition}

The matrix \m{A_1^{\tp}\pi(\gamma)A_0} is square, indexed by \m{I'\times I'}. If the two stabilisers in~(d) carry
little of the mass of~$\omega$, the proposition says that the set where this matrix is singular carries at most
about the share \m{2/\#\F} of the mass.

\begin{proof}
Fix a non-trivial character~$\Psi$ of the additive group of~$\F$; its values have absolute value~$1$. For a finite
set~$I$ and $\rho,\rho'$ in~$\F^I$ write \m{\rho\cdot\rho'=\sum_{i\in I}\rho_i\rho'_i}. The sum of a non-trivial
character over a finite group is zero, and multiplication by a non-zero element of~$\F$ is a bijection of~$\F$.
Hence, coordinate by coordinate,
\begin{equation}\label{e:orth}
\sum_{\rho'\in\F^I}\Psi(\rho\cdot\rho')\ =\ (\#\F)^{\#I}\ \text{ if }\rho=0,\qquad\text{and}\qquad=0\ \text{ otherwise.}
\end{equation}
For $\xi_1,\xi_0\in\F^{I'}$ and $\gamma\in G$ put
\m{\omega_{\xi_1,\xi_0}(\gamma)=\Psi\bigl(\xi_1\cdot A_1^{\tp}\pi(\gamma)A_0\,\xi_0\bigr)\,\omega(\gamma)}.

\emph{Step 1.} Fix $\gamma$ and~$\xi_1$. By~\eqref{e:orth} for the set~$I'$, the mean over~$\xi_0$ of the numbers
\m{\Psi(\xi_1\cdot A_1^{\tp}\pi(\gamma)A_0\,\xi_0)} is $1$ if \m{\xi_1^{\tp}A_1^{\tp}\pi(\gamma)A_0=0} and $0$ otherwise.
The $\xi_1$ with this property form a subspace of~$\F^{I'}$; it has at least one element, and at least $\#\F$
elements when the matrix is singular. So the sum over~$\xi_1$ of the mean over~$\xi_0$ of
\m{\sum_{\gamma}\omega_{\xi_1,\xi_0}(\gamma)} is the sum over~$\gamma$ of $\omega(\gamma)$ times the number of
elements of this subspace. It is a real number, and as \m{\omega\ge0} by~(a),
\[\#\F\cdot\sum_{\substack{\gamma\in G\\ \det(A_1^{\tp}\pi(\gamma)A_0)=0}}\omega(\gamma)\ \le\
\sum_{\xi_1\in\F^{I'}}\ (\#\F)^{-\#I'}\sum_{\xi_0\in\F^{I'}}\ \re\sum_{\gamma\in G}\omega_{\xi_1,\xi_0}(\gamma).\]

\emph{Step 2.} If $\xi_1$ or $\xi_0$ is zero, then \m{\omega_{\xi_1,\xi_0}=\omega}. We claim that
\begin{equation}\label{e:step2}
\re\sum_{\gamma\in G}\omega_{\xi_1,\xi_0}(\gamma)\ \le\ \sqrt{(\#\F)^{\#I}}\cdot Y\qquad\text{for all non-zero }\xi_1,\xi_0\in\F^{I'} .
\end{equation}
In the inner product
space $\CC^{\F^I}$, with standard basis $(\one_\rho)_{\rho\in\F^I}$, put
\[\Psi_\rho\ =\ (\#\F)^{-\#I/2}\sum_{\rho'\in\F^I}\Psi(-\rho\cdot\rho')\,\one_{\rho'}\qquad\text{for }\rho\in\F^I .\]
By~\eqref{e:orth} these vectors form a second orthonormal family, and
\[\ip{\Psi_\rho}{\one_{\rho'}}=(\#\F)^{-\#I/2}\,\Psi(\rho\cdot\rho')\qquad\text{for all }\rho,\rho'\in\F^I .\]
Consider in \m{\CC^{\F^I}\otimes\Hs} the two vectors
\[F=\sum_{\gamma\in G}\Psi_{(\pi(\gamma)^{\tp})^{-1}A_1\xi_1}\otimes v(\gamma),\qquad
F'=\sum_{\gamma\in G}\one_{\pi(\gamma)A_0\xi_0}\otimes v(\gamma).\]
The dot product of \m{(\pi(\gamma)^{\tp})^{-1}A_1\xi_1} and \m{\pi(\gamma')A_0\xi_0} is
\m{\xi_1\cdot A_1^{\tp}\pi(\gamma^{-1}\gamma')A_0\,\xi_0}. So by~(b) the inner product of the term of~$F$ at~$\gamma$
with the term of~$F'$ at~$\gamma'$ is \m{(\#\F)^{-\#I/2}\,\omega_{\xi_1,\xi_0}(\gamma^{-1}\gamma')}, and
Lemma~\ref{l:groupsum} gives
\[\ip{F}{F'}\ =\ \#G\cdot(\#\F)^{-\#I/2}\sum_{\gamma\in G}\omega_{\xi_1,\xi_0}(\gamma).\]
In the same way, the inner product of the terms of~$F$ at $\gamma$ and~$\gamma'$ is $\omega(\gamma^{-1}\gamma')$ if
\m{(\pi(\gamma^{-1}\gamma')^{\tp})^{-1}A_1\xi_1=A_1\xi_1} and $0$ otherwise, because the family $(\Psi_\rho)$ is
orthonormal and \m{\gamma\mapsto(\pi(\gamma)^{\tp})^{-1}} is a homomorphism. By Lemma~\ref{l:groupsum} and~(d),
\m{\norm{F}^2\le\#G\cdot Y}; likewise \m{\norm{F'}^2\le\#G\cdot Y}. Hence
\m{\re\ip{F}{F'}\le\norm{F}\,\norm{F'}\le\tfrac12(\norm{F}^2+\norm{F'}^2)\le\#G\cdot Y}, which is~\eqref{e:step2}.

\emph{Step 3.} In the right side of Step~1, the terms with \m{\xi_1=0} contribute \m{\sum_{\gamma}\omega(\gamma)}. The
terms with \m{\xi_1\ne0} and \m{\xi_0=0} contribute at most the same, because there are fewer than $(\#\F)^{\#I'}$
of them and \m{\sum_{\gamma}\omega(\gamma)\ge0}. The remaining terms contribute at most
\m{(\#\F)^{\#I'}\sqrt{(\#\F)^{\#I}}\cdot Y}, by~\eqref{e:step2} and~(c). This is the assertion.
\end{proof}
\leanline{\lean{minors}\sep uses \lean{groupsum}, and from the background library \lean{orth} (the relation
\eqref{e:orth}), \lean{invT} and \lean{transposeGL} (the map \m{\gamma\mapsto(\pi(\gamma)^{\tp})^{-1}})}

\begin{definition}[the dyad sum of a family]\label{d:dyad}
Let $\Hs$ be a complex inner product space with an orthonormal basis $e=(e_{i'})_{i'\in I'}$ indexed by a finite
set~$I'$, and let $v=(v_i)_{i\in I}$ be a family of vectors of~$\Hs$ indexed by a finite set~$I$. For $i\in I$ let
\m{\bar v_i=(\ip{v_i}{e_{i'}})_{i'\in I'}} in~$\CC^{I'}$ be the vector of the complex conjugates of the coordinates
of~$v_i$ in the basis~$e$. The \kw{dyad sum} of~$v$ is
\[\dyad{e}{v}\ =\ \sum_{i\in I}\bar v_i\otimes v_i\quad\text{in }\CC^{I'}\otimes\Hs .\]
\end{definition}
\leanline{\lean{dyad}}

\begin{lemma}[inner products of dyad sums]\label{l:innerdyad}
Let $\Hs$ be a complex inner product space with an orthonormal basis $e=(e_{i'})_{i'\in I'}$ indexed by a finite
set~$I'$, and let $v=(v_i)_{i\in I}$ and $v'=(v'_{i''})_{i''\in I''}$ be families of vectors of~$\Hs$ indexed by finite
sets $I$ and~$I''$. Then
\[\ip{\dyad{e}{v}}{\dyad{e}{v'}}\ =\ \sum_{i\in I}\ \sum_{i''\in I''}\abs{\ip{v'_{i''}}{v_i}}^{2}.\]
\end{lemma}
\begin{proof}
The inner product of \m{\bar v_i\otimes v_i} and \m{\bar v'_{i''}\otimes v'_{i''}} is
\m{\ip{\bar v_i}{\bar v'_{i''}}\cdot\ip{v_i}{v'_{i''}}}, and by Parseval's identity
\[\ip{\bar v_i}{\bar v'_{i''}}\ =\ \sum_{i'\in I'}\ip{v'_{i''}}{e_{i'}}\,\ip{e_{i'}}{v_i}\ =\ \ip{v'_{i''}}{v_i} .\]
The product of $\ip{v'_{i''}}{v_i}$ and its complex conjugate $\ip{v_i}{v'_{i''}}$ is
$\abs{\ip{v'_{i''}}{v_i}}^2$.
\end{proof}
\leanline{\lean{innerdyad}}

Under the identification of \m{\CC^{I'}\otimes\Hs} with the endomorphisms of~$\Hs$ that the basis~$e$ provides, the
dyad sum of~$v$ is the positive operator \m{F\mapsto\sum_{i\in I}\ip{v_i}{F}\,v_i}, and the inner product of the
lemma is the trace of the product of two such operators. This remark is ours and is not used.

\begin{proposition}[a variance inequality]\label{p:variance}
Let $\Hs$ be a finite-dimensional complex inner product space, $e=(e_{i'})_{i'\in I'}$ an orthonormal basis of~$\Hs$
indexed by a finite set~$I'$, and $I$ a finite set. Let \m{T\colon\CC^I\to\Hs} and \m{C\colon\Hs\to\Hs} be
$\CC$-linear maps, \m{\Us\subseteq\CC^I} and \m{\Hs'\subseteq\Hs} linear subspaces, and $R$ a real number. Suppose
that \m{\sum_{i\in I}\norm{(T^{*}T-1)\one_i}\le R}, where $T^{*}$ is the adjoint of~$T$, and that
\m{(T-C\,T)(\Us)\subseteq\Hs'}. Then
\[\dim\Us\ -\ R\ -\ \sum_{i\in I}\norm{C\,T\one_i}^2\ \ \le\ \ \re\ip{\dyad{e}{T\one}}{\dyad{e}{e'}},\]
where $T\one$ is the family $(T\one_i)_{i\in I}$ and $e'$ is the orthonormal basis of~$\Hs'$, indexed by
\m{\{1,\dots,\dim\Hs'\}}, that the library provides.
\end{proposition}

By Lemma~\ref{l:innerdyad} and Parseval's identity in~$\Hs'$ the right side is the real number
\m{\sum_{i\in I}\norm{p_{\Hs'}T\one_i}^2}, the square of the Hilbert--Schmidt norm of \m{p_{\Hs'}T}; it does not depend
on the choice of~$e'$. So the proposition says: if $T$ is nearly an isometry, and if on a subspace~$\Us$ it lands
in~$\Hs'$ up to the small map~$C\,T$, then the part of~$T$ in~$\Hs'$ has Hilbert--Schmidt norm squared at least about
$\dim\Us$.

\begin{proof}
Let $v=(v_{i''})_{i''\in I''}$ be an orthonormal basis of~$\Us$, so that \m{\#I''=\dim\Us}.

\emph{Step 1.} The endomorphism \m{p_{\Us}(1-T^{*}T)} of~$\CC^I$ maps into~$\Us$, so its trace can be computed in the
basis~$v$ of~$\Us$ and also in the standard basis:
\[\dim\Us-\sum_{i''\in I''}\norm{Tv_{i''}}^2\ =\ \sum_{i''\in I''}\ip{v_{i''}}{(1-T^{*}T)v_{i''}}\ =\ \sum_{i\in I}\ip{p_{\Us}\one_i}{(1-T^{*}T)\one_i}.\]
The left side is real, and the real part of the right side is at most
\[\sum_{i\in I}\norm{p_{\Us}\one_i}\,\norm{(1-T^{*}T)\one_i}\ \le\ R,\]
because \m{\norm{p_{\Us}\one_i}\le1}. Hence
\m{\dim\Us-R\le\sum_{i''\in I''}\norm{Tv_{i''}}^2}.

\emph{Step 2.} For each~$i''$ the vector \m{Tv_{i''}-C\,Tv_{i''}} lies in~$\Hs'$, so $Tv_{i''}$ and $C\,Tv_{i''}$ have the
same component orthogonal to~$\Hs'$. By Pythagoras' theorem,
\[\norm{Tv_{i''}}^2\ \le\ \norm{p_{\Hs'}Tv_{i''}}^2+\norm{C\,Tv_{i''}}^2 .\]

\emph{Step 3.} For every $\CC$-linear map \m{C'\colon\Hs\to\Hs},
\[\sum_{i''\in I''}\norm{C'Tv_{i''}}^2=\sum_{i'\in I'}\ \sum_{i''\in I''}\abs{\ip{(C'T)^{*}e_{i'}}{v_{i''}}}^2
\ \le\ \sum_{i'\in I'}\norm{(C'T)^{*}e_{i'}}^2=\sum_{i\in I}\norm{C'T\one_i}^2,\]
by Parseval's identity in~$\Hs$, Bessel's inequality for the orthonormal family~$v$, and Parseval's identity
in~$\CC^I$.

Steps 1 and~2, and Step~3 for \m{C'=p_{\Hs'}} and for \m{C'=C}, give
\[\dim\Us-R\ \le\ \sum_{i\in I}\norm{p_{\Hs'}T\one_i}^2+\sum_{i\in I}\norm{C\,T\one_i}^2 ,\]
and the first sum is the right side of the assertion, as was said after the statement.
\end{proof}
\leanline{\lean{variance}\sep uses \lean{innerdyad}}

\sect{Where the results are used}{s:use}
The heights serve Theorem~C of the argument, and the linear algebra of §§\ref{s:few} and~\ref{s:moments} serves
Theorem~B. The notation of this section is that of the overview of the handout: $D$ is admissible data, with third
scale~$b(D)$ and number of lifts~$m(D)$; $\Row(D)$ is the space of rows; for an $m(D)$-fold twist~$\vartheta$ of~$D$,
\m{\Theta_{D,\vartheta}\colon\Row(D)\to\mathcal{A}_{D,\vartheta}} is the chain of lifts and
\m{\mathfrak{z}_D(\vartheta)=\det(\Theta_{D,\vartheta}^{*}\Theta_{D,\vartheta})}. What follows was read in the
statements and in the last steps of the formal proofs; the chapters named give the proofs.

\subhead{Theorem C: a covolume is bounded from below.} Let \m{y_1,\dots,y_n} be linearly independent rows with
\m{n\ge(1-4/b(D)^2)\dim\Row(D)}, of $D$-index at most~$\delta$, whose lifts at the trivial twist satisfy
\m{\covol_\sigma^2(\Theta_{D,1}y_1,\dots,\Theta_{D,1}y_n)\le\exp(V)} for every $\sigma\in\AutC$. Chapter~VIII provides
a family $\theta=(\theta_{\ell,\iota})$ of $m(D)$-fold twists, one for each pair in~$\IX$, that is fit for~$D$. Two
height bounds are then proved.

\emph{First bound.} The family \m{\mathfrak{z}_D(\theta_{\ell,\iota})/\mathfrak{z}_D(1)} has height at most
\[V+2\delta+b(D)\cdot\dim\Row(D).\]
This is Theorem~2 for \m{W=\Row(D)}, for the set~$\Xi$ of the $m(D)$-fold twists with
\m{\vartheta_0=1}, for \m{T_\vartheta=\Theta_{D,\vartheta}}, and for the sets $\Ls_{\ell,\iota}(j)$ of the
$j$-integral forms, for which (f) is the integrality of the pairing. The basis~$v'$ consists of \m{y_1,\dots,y_n}
and of~$n'$ further rows \m{z_1,\dots,z_{n'}} given by Lemma~\ref{l:extend}, taken from a spanning set of crude
rows that Chapter~VII constructs. These come with scalars~$c_{\ell,\iota}$, ratios of a twisted to an untwisted constant, for which the family
$\abs{c}^{-2}$ has a height bound. The matrix of~(g) is made of two blocks: the matrix that the $D$-index
provides for the~$y_i$, with \m{\abs{j(\det u)}_\ell\ge\exp(-\eta_\ell)}, and $c_{\ell,\iota}$ times the unit matrix
for the~$z_i$; so \m{\nu=n'}, and hypothesis~(e) holds with the $\delta$ of the $D$-index. The number~$\alpha$ is $V$
plus a term for the~$z_i$. There are at most \m{(4/b(D)^2)\dim\Row(D)} vectors~$z_i$, and the last lines of the formal
proof bound what each of them adds to $\alpha$ and to~$\nu h$ by~$b(D)^3/4$: this is the term \m{b(D)\cdot\dim\Row(D)}.

\emph{Second bound.} The family \m{\mathfrak{z}_D(1)^2/\mathfrak{z}_D(\theta_{\ell,\iota})^2} has height at most
\[2\,\dim\Row(D)\cdot\bigl(2\log b(D)\cdot\budget(D)-\hmain(D)\bigr).\]
Here $\mathfrak{z}_D(\vartheta)$ is a constant times the
product of the volumes of the $m(D)$ single lifts, each ratio of an untwisted to a twisted volume has a height
bound (Chapter~VII), and Proposition~\ref{p:products} and Lemma~\ref{l:mono} combine them.

\begin{center}
\begin{tikzpicture}[x=1mm,y=1mm,>={Stealth[length=1.6mm]},
  bx/.style={draw=rulegrey,line width=0.5pt,rounded corners=1.2pt,fill=arith!3!white,inner xsep=5pt,inner ysep=4pt,
    align=center,font=\fontsize{8.3}{10.4}\selectfont},
  gx/.style={bx,fill=inksoft!5!white},
  lb/.style={font=\fontsize{7.2}{8.8}\selectfont\color{inksoft},align=center}]
\node[gx] (a) at (0,9) {first bound:\\index and size (Theorem~2)};
\node[gx] (b) at (0,-9) {second bound:\\congruences (Chapters VII, VIII)};
\node[bx] (c) at (52,0) {the square of the first family\\times the second is~$1$};
\node[bx] (d) at (101,0) {Theorem~C};
\draw[->,inksoft] (a.east) -- (c);
\draw[->,inksoft] (b.east) -- (c);
\draw[->,inksoft] (c) -- node[lb,above=1pt] {Theorem~1} (d);
\end{tikzpicture}
\end{center}

The square of the first family times the second is the constant family~$1$. By
Proposition~\ref{p:products} and Theorem~1, twice the first bound plus the second is non-negative:
\[\dim\Row(D)\cdot\bigl(\hmain(D)-2\log b(D)\cdot\budget(D)-b(D)\bigr)\ \le\ V+2\delta,\]
which gives Theorem~C because \m{b(D)\le\log b(D)\cdot\budget(D)}. The twist $\theta_{\ell,\iota}$ depends on the pair
$(\ell,\iota)$, so the numbers $\mathfrak{z}_D(\theta_{\ell,\iota})$ are not the conjugates of one algebraic number,
and no product formula applies to them. The height carries the $\ell$-adic bounds of the two families through the
multiplication, and the product formula is applied at the end to the auxiliary numbers~$H$ alone. The complete
deduction belongs to the note on the unit \emph{New.Endgame}, which ends the argument.

\subhead{Other height bounds.} Height bounds for families of ratios of a twisted to an untwisted quantity
(values of Hecke $L$-functions, volumes of single lifts, the constants above) are proved in the units
\emph{New.Rows}, \emph{New.Exchange} and \emph{New.Selmer}, and combined by Lemma~\ref{l:mono} and
Proposition~\ref{p:products}.

\subhead{Theorem B: enough forms.} Theorem~3 is applied with \m{W=\Row(D)}, with $U$ the space of columns
$\Col(D)$, with the finite group~$\mathsf{G}(D)$ as index set and with \m{\varepsilon=b(D)^{-2}}. The map~$B_\gamma$ is
a fixed operator from the columns onto the ramified columns followed by the Poincaré map
\m{\mathcal{P}_D(\gamma)}, in the normalisation of Chapter~VI. That this family is $b(D)^{-2}$-spread is the filling
theorem of that chapter. Its proof is a second moment over the group: Proposition~\ref{p:variance} is applied to
each translate and summed, with Lemmas \ref{l:groupsum} and~\ref{l:innerdyad}, and Proposition~\ref{p:minors} is
applied inside an induction, to a function of positive type and to the set where a minor is singular. Theorem~3
then gives the finite set~$\Ts$ of Theorem~B, with \m{\#\Ts\cdot\dim\Col(D)\le5\,\dim\Row(D)\cdot\log(1+b(D)^2)} and a
span of dimension at least \m{(1-4/b(D)^2)\dim\Row(D)}.

\sect{Notes on the formal text}{s:notes}
\begin{enumerate}[label=\textup{\arabic*.},leftmargin=1.6em,itemsep=5pt]
\item \emph{Two ways to add up the local exponents.} Definition~\ref{d:height} asks for the bound on the sums
of~$\eta$ over the primes below~$n$, for all large~$n$. Lemma~\ref{l:ofle} assumes the bound for all finite sets
of primes, and hypothesis~(e) of Theorem~2 for all finite sets of natural numbers. The formal text has these three
forms, and they are kept here. In the definition the formal text says that the bound holds eventually; the
threshold~$n_0$ is this statement written out.
\item \emph{Functions defined everywhere.} The formal Lemma~\ref{l:pf} includes \m{q=0}, by the convention
\m{\log0=0}. Definition~\ref{d:spread} is stated for arbitrary complex vector spaces, with the library's dimension;
Theorem~3 uses it for finite-dimensional spaces only.
\item \emph{Theorem 2.} The formal statement asks (f) and~(g) for every homomorphism \m{j\colon\CC\to\Cl}, and its
proof uses two of them for each $(\ell,\iota)$. The sets $\Ls_{\ell,\iota}(j)$ are arbitrary subsets. The formal
conclusion is written with the volume factor of the library, squared, and the quotient is formed in~$\CC$. The cost
$2\delta$ comes from the two matrices of the proof, one for $\jof{\iota}$ and one for its conjugate, and
$\nu h$ from the family $\abs{c}^{-2}$, which is bounded independently. The base map $T_{\vartheta_0}$ must be
injective, so that the witness $\Gamma_{\vartheta_0}$ is not zero.
\item \emph{Theorem 3.} The formal statement takes $W$ to be an inner product space; the proof does not use the
inner product. It chooses an admissible set whose span is maximal for inclusion; the proof above takes the
largest dimension, which amounts to the same.
\item \emph{Proposition~\ref{p:minors}.} The formal statement gives the pair~$A$ as a function on the two truth
values; $A_1$ is its value at ``true'' and $A_0$ its value at ``false''. The condition in~(d) is written with one
formula for both, through an automorphism of $\GL_I(\F)$ that is the identity for $A_0$ and the passage to the
inverse of the transpose for~$A_1$, a definition of the background library. The formal proof takes a character
from Mathlib and the relation~\eqref{e:orth} from the background library, where it is stated for finite commutative
rings; here it is proved for a field in one line.
\item \emph{Proposition~\ref{p:variance}.} The basis~$e$ enters the statement only through the dyad sums, and the
basis~$e'$ is the one that the library attaches to a finite-dimensional inner product space.
\item \emph{Arguments read from the text.} The list of arguments of each declaration was made by a program from
the text of the files. For some declarations of the second file it contains variables of the surrounding
section whose names are bound again inside the statement (a basis, a second map on~$G$, a field, an index set).
Lean does not take them as arguments, as the dossier notes, and where such a variable is explicit the uses of the
lemma in the artifact confirm it. In the accounts of this note they are matched with the bound variables of the
same name.
\item \emph{What the final theorems use.} All declarations of the unit are stated here. Whether each of them is
reached from the final theorems is for the kernel-based ledger of the handout to say.
\item \emph{Literature named by the formal text.} The file on heights calls the product formula through a norm
and adelic norms changed place by place standard, and names as related work Gaudron (2006), Bombieri and Gubler
(2006, ch.~1) and notes of K.~Conrad. The file on determinants names Wolsey (1982) for greedy covering with a
logarithmic loss, and Bost (2001) and Laurent (1994) for size against divisibility of a determinant. The
references are listed below as the artifact gives them, with the state that its own list of references records;
we have not compared them.
\end{enumerate}

\subhead{How this note was written.} From the dossier of the unit: its statements in words, its structured
proofs, and its imports. The program check of the dossier covered \UnitSteps{} leaf steps of \UnitProofs{}
proofs, found the \UnitMembers{} statements at their lines and compared the quotations of the imports with the
source; a second reading of the dossier by the reading agent asked for no correction. I read the files of the
unit completely, statements and proofs, and compared every statement of this note with its Lean text argument by
argument, in Lean's order of quantifiers. I also read: the statements of the imported declarations in the
background library and the proof of Input~1; the formal proofs of Theorem~C and of its two height bounds, the
statement and the last lines of the proof of Theorem~B, the statement of the filling theorem, and the places where
the two inequalities of §\ref{s:moments} are applied; the authors' list of references. I did not read Mathlib, and
I did not run a build. Shorter than, or different from, the formal proofs: Lemma~\ref{l:pf} (the formal proof is
one call of automation on the prime factorisation); the bound for $\norm{F}^2$ in Proposition~\ref{p:minors}, which
the formal proof obtains from one general identity for an orthonormal family and an automorphism of the linear
group; Step~1 of Proposition~\ref{p:variance}, which follows the same two traces; the arithmetic at the end of
Theorem~3, which the formal proof leaves to automation and which is written out here. Lemma~\ref{l:gram},
Corollary~\ref{c:reciprocal}, Remark~\ref{r:constant}, the examples and the diagram are mine.

\head{Appendix A. Concordance with the Lean artifact}
{\fontsize{8.6}{11.6}\selectfont\renewcommand{\leanfont}{\ttfamily\fontsize{7.7}{10}\selectfont\color{leancol}}
\setlength{\LTleft}{0pt}\setlength{\LTright}{\fill}\setlength{\tabcolsep}{0pt}
\begin{longtable}{@{}p{30mm}p{56mm}p{43mm}r@{}}
\textbf{In this note} & \textbf{Lean declaration} & \textbf{File} & \textbf{Lines}\\\midrule\endhead
\ConcordanceRows
\end{longtable}}

\subhead{Letters of the formal text and of this note.} The letters of this note follow the overview of the
handout where it has fixed them, and differ from those of the formal text. Heights: the formal text writes
\lf{c}, \lf{y} for our $h$, $Z$; \lf{e} for~$\eta$; \lf{M} for the bound~$n$ of the primes. Section~\ref{s:moments}:
it writes \lf{H} for~$\Hs$; \lf{k} for~$\F$; \lf{n}, \lf{m}, \lf{l} for the index sets $I$, $I'$, $I''$; \lf{u},
\lf{v} for the maps $v$, $v'$ on~$G$; \lf{a}, \lf{b}, \lf{s} for the elements of~$G$; \lf{f} and \lf{Θ}
for~$\omega$; \lf{b} for the basis~$e$; \lf{P}, \lf{R}, \lf{U}, \lf{A}, \lf{ε} for $T$, $C$, $\Us$, $\Hs'$, $R$.
Theorem~3: \lf{V}, \lf{W} for $U$, $W$; \lf{Θ} or \lf{S} for~$\Omega$; \lf{η} for~$\varepsilon$; \lf{F} for~$W'$;
\lf{t} for~$\Ts$. Lemma~\ref{l:extend}: \lf{V} for~$W$ and \lf{k} for~$n'$. Theorem~2: \lf{Θ} for~$\Xi$;
\lf{H ϑ} for~$\Hs_\vartheta$; \lf{u} for the maps~$T$; \lf{o} for~$\vartheta_0$; \lf{S} and \lf{x} for $I$
and~$v'$; \lf{g} for~$h$; \lf{γ} for~$c$; \lf{e} for~$\eta$; \lf{U} for~$J$; \lf{Λ} for~$\Ls$; \lf{k} for~$\nu$;
\lf{m} for the matrix~$u$.

\subhead{Statements as they are written in the artifact.} The definition of a height bound and Theorem~1:
\quoted{hle}
\quoted{nonneg}
Theorem~2:
\quoted{onedet}
The spread property and Theorem~3:
\quoted{spread}
\quoted{few}
Propositions \ref{p:minors} and~\ref{p:variance}:
\quoted{minors}
\quoted{variance}

\head{Notation}
\begingroup\fontsize{8.6}{10.9}\selectfont\raggedright
\setlength{\columnsep}{9mm}\setlength{\parskip}{0pt}
\begin{multicols}{2}
\nt{$\ell$}{a prime number}
\nt{$\Ql$, $\Qlb$, $\Cl$}{the $\ell$-adic numbers, an algebraic closure, its completion}
\nt{$\abs{\,\cdot\,}_\ell$, $\abs{\,\cdot\,}$}{the absolute value with $\abs{\ell}_\ell=1/\ell$; the absolute value of~$\CC$}
\nt{$\iota$}{an isomorphism of rings from $\CC$ to $\Qlb$, or to a field~$L$}
\nt{$\IX$}{the set of the pairs $(\ell,\iota)$}
\nt{$\sigma$, $\AutC$}{a field automorphism of~$\CC$; the group of all of them}
\nt{$Z$, $Z'$, $Z_i$}{families of complex numbers indexed by~$\IX$, with entries $Z_{\ell,\iota}$}
\nt{$h$, $h'$, $h_i$}{real numbers, bounds for heights}
\nt{$\hle{h}{Z}$}{$Z$ has height at most~$h$}
\nt{$H$, $\alpha$, $\eta$}{a witness: a non-zero complex number, a real number, a sequence $(\eta_0,\eta_1,\dots)$ of real numbers}
\nt{$n$, $n_0$, $n'$}{natural numbers; a threshold; the number of added vectors}
\nt{$q$, $q_{\mathrm{num}}$, $q_{\mathrm{den}}$}{a rational number, its numerator and denominator}
\nt{$\ord_\ell$, $\Nm$, $\deg H$}{exponent of~$\ell$; norm of a field extension; degree of $\QQ(H)$ over~$\QQ$}
\nt{$\bar h$}{in Remark~\ref{r:constant}: a mean of logarithms}
\nt{$j$, $\bar j$, $\jof{\iota}$}{a homomorphism of rings $\CC\to\Cl$; $j$ after complex conjugation; $\iota$ followed by the inclusion}
\nt{$O_j$}{the complex numbers with $\abs{j(\,\cdot\,)}_\ell\le1$}
\nt{$L$; $\tau$, $\tau_0$}{an algebraically closed field; injective homomorphisms of rings}
\nt{$Y$}{a real number, a bound}
\nt{$I$, $I'$, $I''$; $i$, $i'$, $i''$, $i_1$}{finite sets; their elements (in Theorem~2 and Lemma~\ref{l:gram}, $i$ and $i'$ both run over~$I$)}
\nt{$\#I$}{the number of elements}
\nt{$\Hs$, $\Hs'$, $\Hs_\vartheta$}{complex inner product spaces; $\Hs'$ a subspace of~$\Hs$}
\nt{$\ip{\,\cdot\,}{\,\cdot\,}$, $\norm{\,\cdot\,}$}{inner product, conjugate-linear in the first variable; its norm}
\nt{$F$, $F'$, $F''$, $F_i$}{vectors of an inner product space}
\nt{$v$, $v'$}{families of vectors, or maps from~$G$ to~$\Hs$; $Tv$ the family $(Tv_i)$}
\nt{$\Gram(v)$}{the Gram matrix $(\ip{v_i}{v_{i'}})$}
\nt{$T$, $T_\vartheta$, $T^{*}$}{$\CC$-linear maps into inner product spaces; the adjoint; $T^{*}T$ as in §\ref{s:intro}}
\nt{$W$, $W'$, $W_y$, $W_{\Ts}$}{a finite-dimensional complex vector space; subspaces}
\nt{$\Xi$, $\vartheta_0$, $\theta$}{a set; one of its elements; a family $(\theta_{\ell,\iota})$ of elements}
\nt{$\vartheta$}{In Sections~\ref{s:intro} to~\ref{s:moments} also an element of the set $\Xi$, and in Section~\ref{s:use} a twist of the data (Convention~0.5)}
\nt{$c$}{a family of non-zero complex numbers $c_{\ell,\iota}$}
\nt{$\delta$}{a real number, a bound for an index}
\nt{$\nu$}{In Sections~\ref{s:intro} to~\ref{s:use} also a natural number, never a character of the ideles (Convention~0.6)}
\nt{$J$}{a finite set of natural numbers}
\nt{$\Ls_{\ell,\iota}(j)$}{a subset of $\Hs_{\theta_{\ell,\iota}}$}
\nt{$u$, $u'$, $\overline{u'}$, $u^{\tp}$}{matrices of complex numbers; conjugate entries; transpose}
\nt{$\Gamma_\vartheta$}{$\det\Gram(T_\vartheta v')$}
\nt{$U$, $\Omega$, $B$, $B_\gamma$}{a complex vector space; a set; a family of linear maps $U\to W$}
\nt{$\varepsilon$, $\spread{B}{\varepsilon}$}{a real number; $B$ is $\varepsilon$-spread}
\nt{$\Ts$, $\Ts'$, $\Ts''$}{finite subsets of~$\Omega$}
\nt{$\Pi$, $\mathrm{rk}$, $\codim$}{the potential, a rank, a codimension, in the proof of Theorem~3}
\nt{$S$; $y_i$, $z_i$}{a spanning set; vectors}
\nt{$G$, $\gamma$, $\gamma'$}{a finite group; elements of~$G$, or of~$\Omega$}
\nt{$\omega$, $\omega_{\xi_1,\xi_0}$}{a function on~$G$; its product with character values}
\nt{$\F$, $\GL_I(\F)$, $\pi$}{a finite field; invertible matrices; a homomorphism $G\to\GL_I(\F)$}
\nt{$A=(A_1,A_0)$}{a pair of matrices over~$\F$}
\nt{$\xi$, $\xi'$, $\xi_1$, $\xi_0$, $\rho$, $\rho'$}{vectors with coordinates indexed by a finite set}
\nt{$\Psi$, $\Psi_\rho$}{a non-trivial additive character of~$\F$; the normalised characters of~$\F^I$}
\nt{$\one$, $\one_i$}{the standard basis of $\CC^I$}
\nt{$e$, $e'$}{orthonormal bases of $\Hs$ and of~$\Hs'$}
\nt{$\dyad{e}{v}$, $\bar v_i$}{the dyad sum; conjugate coordinates}
\nt{$C$, $C'$, $\Us$, $R$}{linear maps $\Hs\to\Hs$; a subspace of~$\CC^I$; a real number, a bound}
\nt{$p_{\Us}$, $p_{\Hs'}$}{orthogonal projections}
\nt{$D$, $b(D)$, $m(D)$, $V$}{in §\ref{s:use}: admissible data, its third scale, its number of lifts; a real number}
\nt{$\Row(D)$, $\Col(D)$, $\mathsf{G}(D)$}{rows, columns, the finite group of the data}
\nt{$\Theta_{D,\vartheta}$, $\mathcal{A}_{D,\vartheta}$, $\mathfrak{z}_D(\vartheta)$}{the chain of lifts, its target, $\det(\Theta_{D,\vartheta}^{*}\Theta_{D,\vartheta})$}
\nt{$\mathcal{P}_D(\gamma)$}{the Poincaré map translated by~$\gamma$}
\nt{$\covol_\sigma$}{in §\ref{s:use}: the covolume at~$\sigma$, as in the overview of the handout}
\nt{$\budget(D)$, $\hmain(D)$}{the budget and the main amount of the data}
\end{multicols}
\endgroup

\head{References}
{\fontsize{9.2}{12.6}\selectfont
\noindent\hangindent=2.2em\hangafter=1 [St]\enspace E.~Steinitz, Algebraische Theorie der Körper, J.~reine angew.\
Math.\ 137 (1910) 167--309. Taken from the catalogue of quoted results of this handout; place not confirmed.\par
\vspace{4pt}\noindent Named by the formal text as related work, copied from its list of references with the state recorded
there; not compared by us. For the file on heights: Gaudron, Pentes des fibrés vectoriels adéliques sur un corps
global, 2006 (work found, no place given); Bombieri--Gubler, Heights in Diophantine Geometry, 2006 (place confirmed
at second hand); K.~Conrad, Zorn's lemma and some applications~II, notes (not compared, no place given). For the
file on determinants: Wolsey, Greedy algorithm for the submodular set covering problem, 1982; Bost, Algebraic
leaves of algebraic foliations over number fields, 2001; Laurent, Linear forms in two logarithms and interpolation
determinants, 1994 (for these three: work found, no place given).\par}

\end{document}
